% Options for packages loaded elsewhere
\PassOptionsToPackage{unicode}{hyperref}
\PassOptionsToPackage{hyphens}{url}
\documentclass[
  12pt,
]{article}
\usepackage{xcolor}
\usepackage[margin=1in]{geometry}
\usepackage{amsmath,amssymb}
\setcounter{secnumdepth}{-\maxdimen} % remove section numbering
\usepackage{iftex}
\ifPDFTeX
  \usepackage[T1]{fontenc}
  \usepackage[utf8]{inputenc}
  \usepackage{textcomp} % provide euro and other symbols
\else % if luatex or xetex
  \usepackage{unicode-math} % this also loads fontspec
  \defaultfontfeatures{Scale=MatchLowercase}
  \defaultfontfeatures[\rmfamily]{Ligatures=TeX,Scale=1}
\fi
\usepackage{lmodern}
\ifPDFTeX\else
  % xetex/luatex font selection
\fi
% Use upquote if available, for straight quotes in verbatim environments
\IfFileExists{upquote.sty}{\usepackage{upquote}}{}
\IfFileExists{microtype.sty}{% use microtype if available
  \usepackage[]{microtype}
  \UseMicrotypeSet[protrusion]{basicmath} % disable protrusion for tt fonts
}{}
\usepackage{setspace}
\makeatletter
\@ifundefined{KOMAClassName}{% if non-KOMA class
  \IfFileExists{parskip.sty}{%
    \usepackage{parskip}
  }{% else
    \setlength{\parindent}{0pt}
    \setlength{\parskip}{6pt plus 2pt minus 1pt}}
}{% if KOMA class
  \KOMAoptions{parskip=half}}
\makeatother
\usepackage{longtable,booktabs,array}
\newcounter{none} % for unnumbered tables
\usepackage{calc} % for calculating minipage widths
% Correct order of tables after \paragraph or \subparagraph
\usepackage{etoolbox}
\makeatletter
\patchcmd\longtable{\par}{\if@noskipsec\mbox{}\fi\par}{}{}
\makeatother
% Allow footnotes in longtable head/foot
\IfFileExists{footnotehyper.sty}{\usepackage{footnotehyper}}{\usepackage{footnote}}
\makesavenoteenv{longtable}
\usepackage{graphicx}
\makeatletter
\newsavebox\pandoc@box
\newcommand*\pandocbounded[1]{% scales image to fit in text height/width
  \sbox\pandoc@box{#1}%
  \Gscale@div\@tempa{\textheight}{\dimexpr\ht\pandoc@box+\dp\pandoc@box\relax}%
  \Gscale@div\@tempb{\linewidth}{\wd\pandoc@box}%
  \ifdim\@tempb\p@<\@tempa\p@\let\@tempa\@tempb\fi% select the smaller of both
  \ifdim\@tempa\p@<\p@\scalebox{\@tempa}{\usebox\pandoc@box}%
  \else\usebox{\pandoc@box}%
  \fi%
}
% Set default figure placement to htbp
\def\fps@figure{htbp}
\makeatother
% definitions for citeproc citations
\NewDocumentCommand\citeproctext{}{}
\NewDocumentCommand\citeproc{mm}{%
  \begingroup\def\citeproctext{#2}\cite{#1}\endgroup}
\makeatletter
 % allow citations to break across lines
 \let\@cite@ofmt\@firstofone
 % avoid brackets around text for \cite:
 \def\@biblabel#1{}
 \def\@cite#1#2{{#1\if@tempswa , #2\fi}}
\makeatother
\newlength{\cslhangindent}
\setlength{\cslhangindent}{1.5em}
\newlength{\csllabelwidth}
\setlength{\csllabelwidth}{3em}
\newenvironment{CSLReferences}[2] % #1 hanging-indent, #2 entry-spacing
 {\begin{list}{}{%
  \setlength{\itemindent}{0pt}
  \setlength{\leftmargin}{0pt}
  \setlength{\parsep}{0pt}
  % turn on hanging indent if param 1 is 1
  \ifodd #1
   \setlength{\leftmargin}{\cslhangindent}
   \setlength{\itemindent}{-1\cslhangindent}
  \fi
  % set entry spacing
  \setlength{\itemsep}{#2\baselineskip}}}
 {\end{list}}
\usepackage{calc}
\newcommand{\CSLBlock}[1]{\hfill\break\parbox[t]{\linewidth}{\strut\ignorespaces#1\strut}}
\newcommand{\CSLLeftMargin}[1]{\parbox[t]{\csllabelwidth}{\strut#1\strut}}
\newcommand{\CSLRightInline}[1]{\parbox[t]{\linewidth - \csllabelwidth}{\strut#1\strut}}
\newcommand{\CSLIndent}[1]{\hspace{\cslhangindent}#1}
\setlength{\emergencystretch}{3em} % prevent overfull lines
\providecommand{\tightlist}{%
  \setlength{\itemsep}{0pt}\setlength{\parskip}{0pt}}
\usepackage{float}
\usepackage{tabularray}
\usepackage[normalem]{ulem}
\usepackage{graphicx}
\usepackage{rotating}
\UseTblrLibrary{siunitx}
\NewTableCommand{\tinytableDefineColor}[3]{\definecolor{#1}{#2}{#3}}
\newcommand{\tinytableTabularrayUnderline}[1]{\underline{#1}}
\newcommand{\tinytableTabularrayStrikeout}[1]{\sout{#1}}
\usepackage{bookmark}
\IfFileExists{xurl.sty}{\usepackage{xurl}}{} % add URL line breaks if available
\urlstyle{same}
\hypersetup{
  pdftitle={Contrasting dermal microbial genomes within an urban-rural matrix in salamanders},
  pdfauthor={Marques, A.J.D.1,2,3, Gaun, N.4, Velo-Antón, G.5, Alberdi, A.4},
  hidelinks,
  pdfcreator={LaTeX via pandoc}}

\title{Contrasting dermal microbial genomes within an urban-rural matrix
in salamanders}
\author{Marques, A.J.D.\textsuperscript{1,2,3}, Gaun,
N.\textsuperscript{4}, Velo-Antón, G.\textsuperscript{5}, Alberdi,
A.\textsuperscript{4}}
\date{2026-05-27}

\begin{document}
\maketitle

\setstretch{2}
\textsuperscript{1} Departamento de Biologia, Faculdade de Ciências,
Universidade do Porto, 4099-002 Porto, Portugal

\textsuperscript{2} CIBIO, Centro de Investigação em Biodiversidade e
Recursos Genéticos, InBIO Laboratório Associado, Campus de Vairão,
Universidade do Porto, 4485-661 Vairão, Portugal

\textsuperscript{3} BIOPOLIS Program in Genomics, Biodiversity and Land
Planning, CIBIO, Campus de Vairão, 4485-661 Vairão, Portugal

\textsuperscript{4} Center for Evolutionary Hologenomics / GLOBE
Institute, University of Copenhagen, 1353 Copenhagen K, Denmark

⁵ Universidad de Vigo, ECOEVO Lab, EE Forestal, Campus Universitario A
Xunqueira, E-36005, Pontevedra, Spain

AJDM (ORCID: 0000-0002-6070-0506); NG (ORCID: 0009-0002-3068-4653); GVA
(ORCID: 0000-0002-9483-5695); AA (ORCID: 0000-0002-2875-6446)

\newpage

\section{Abstract}\label{abstract}

The dermis is an essential organ for respiration, hydration, and defence
among amphibians, and its health is at least in part regulated by the
resident microbiota. The dermal microbiome is also very sensitive to
habitat changes such as urbanization, which introduces many novel
environmental pressures to amphibians. There are limited examples,
however, on how urbanization alters amphibian microbiomes. We therefore
aimed to characterize the fire salamander (\emph{Salamandra salamandra}
Linnaeus, 1758) dermal microbiome and examine how it varies in
composition and functionality between urban and rural habitats. We
utilized genome-resolved metagenomics to assemble microbial genomes,
profile microbial communities, and estimate their functional
capabilities. Microbiome diversity estimates within either habitat were
comparable, however, differences were observed in microbial composition.
Relative abundances of anti-fungal and possibly pathogenic microbiota
were found to be greater among rural salamanders. Significant
differences in the inferred functionality of microbiome assemblages were
also identified and warrants further investigation into the
environmental conditions that drive these changes and the implications
for the host dermal health.

\section{Introduction}\label{introduction}

Considered one of the earliest adaptations among vertebrates to
terrestrial life, amphibians have a thin, permeable skin that allows for
gas and water exchange (\citeproc{ref-akat_comden_complex_2023}{Akat
Çömden et al. 2023}). Amphibians' skin comprises most of their
respiratory capillaries and almost completely facilitates respiration
(\citeproc{ref-tattersall_skin_2007}{Tattersall 2007}). This increased
permeability, however, comes at the expense of increased water loss
(\citeproc{ref-toledo_cutaneous_1993}{Toledo and Jared 1993}) and
vulnerability to xenobiotics
(\citeproc{ref-kaufmann_adaption_2016}{Kaufmann and Dohmen 2016};
\citeproc{ref-quaranta_why_2009}{Quaranta et al. 2009}) and pollutants
(\citeproc{ref-purucker_estimating_2023}{Purucker et al. 2023}). The
sensitivity of amphibians to environmental stressors has led to the
generalization of amphibians as indicators of biodiversity loss
(\citeproc{ref-morell_frogs_1999}{Morell 1999};
\citeproc{ref-norris_ghosts_2007}{Norris 2007}), and constitute the most
threatened vertebrate group on Earth, with 41\% of species classified as
globally threatened (\citeproc{ref-cox_global_2022}{Cox et al. 2022};
\citeproc{ref-luedtke_ongoing_2023}{{Luedtke et al.} 2023}). Given the
importance of amphibians' dermis to individual health, much attention
has been directed toward understanding the role that dermal microbiota
have in maintaining dermis health and function
(\citeproc{ref-hernandez-gomez_organismal_2023}{Hernández-Gómez and Hua
2023}; \citeproc{ref-zhu_skin_2024}{Zhu et al. 2024}). Dermal microbiota
contribute to several functions for amphibian hosts, including chemical
signalling (\citeproc{ref-kueneman_amphibian_2014}{Kueneman et al.
2014}), maintenance of the mucosal layer
(\citeproc{ref-woodhams_adaptive_2023}{Woodhams et al. 2023}) which
retains moisture and protect the dermis
(\citeproc{ref-demori_peptides_2019}{{Demori et al.} 2019};
\citeproc{ref-pereira_skin_2018}{Pereira et al. 2018}), and producing
bacterial metabolites contributing to the hosts' defence against
pathogens (\citeproc{ref-muletz-wolz_inhibition_2017}{Muletz-Wolz et al.
2017}; \citeproc{ref-le_sage_preparatory_2021}{Le Sage et al. 2021};
\citeproc{ref-gonzalez-serrano_amphibian_2025}{{González-Serrano et al.}
2025}). In the wake of catastrophic amphibian mortality driven by
chytrid fungal pathogens, most notably \emph{Batrachochytrium
dendrobatidis} (\emph{Bd}; Longcore et al.
(\citeproc{ref-longcore_batrachochytrium_1999}{1999})) and
\emph{Batrachochytrium salamandrivorans} (\emph{Bsal}; {Martel et al.}
(\citeproc{ref-martel_batrachochytrium_2013}{2013}))
(\citeproc{ref-fisher_chytrid_2020}{Fisher and Garner 2020};
\citeproc{ref-ohanlon_recent_2018}{{O'Hanlon et al.} 2018};
\citeproc{ref-rollins-smith_batrachochytrium_2021}{Rollins-Smith and Le
Sage 2021}; \citeproc{ref-castro_monzon_batrachochytrium_2022}{Castro
Monzon et al. 2022}), contemporary research on amphibian microbiota has
largely focused on their role in antifungal defence
(\citeproc{ref-ross_skin_2019}{Ross et al. 2019}). While recent focus
has turned toward understanding the impact that fungal infections have
on disrupting dermal microbiota and how specific microbial taxa react
post-infection, less is known about how environmental conditions
structure the dermal microbiome prior to infection. Identifying pathogen
resistant microbiota is key, but this must also include understanding
the environmental conditions that determine microbiome composition and
subsequently amphibian populations vulnerability.

The amphibian dermal microbiome varies in composition depending on the
host species, life stage, or in response to environmental disturbances
(\citeproc{ref-kueneman_amphibian_2014}{Kueneman et al. 2014};
\citeproc{ref-jani_amphibian_2021}{Jani et al. 2021};
\citeproc{ref-bates_captivity_2019}{Bates et al. 2019};
\citeproc{ref-santos_skin_2023}{Santos et al. 2023}). Environmental
conditions, such as changes in local plant communities
(\citeproc{ref-hernandez-gomez_invasive_2020}{Hernández-Gómez et al.
2020}; \citeproc{ref-dong_rice_2024}{Dong et al. 2024}) and exposure to
pollution (\citeproc{ref-cordeiro_amphibian_2024}{{Cordeiro et al.}
2024}; \citeproc{ref-su_neglected_2023}{Su et al. 2023};
\citeproc{ref-tornabene_energy-related_2023}{Tornabene et al. 2023}) can
alter microbiome assemblages. Repeatedly, studies have suggested that
alterations to microbial communities increased susceptibility to
pathogens (\citeproc{ref-becker_cutaneous_2010}{Becker and Harris 2010};
\citeproc{ref-bletz_disruption_2018}{{Bletz et al.} 2018}) and
contributed to the formation of lesions
(\citeproc{ref-zhu_skin_2024}{Zhu et al. 2024};
\citeproc{ref-luddecke_salamanders_2018}{Lüddecke et al. 2018}). Many of
the disturbances that are found to negatively impact amphibian fitness
and species richness can be attributed to anthropogenic disturbances
(\citeproc{ref-macdonald_meta-analysis_2023}{Macdonald et al. 2023};
\citeproc{ref-iglesias-carrasco_urban_2017}{Iglesias-Carrasco et al.
2017}), with habitat loss and degradation through agriculture and
urbanization, serving as the primary contributor to amphibian declines
(\citeproc{ref-luedtke_ongoing_2023}{{Luedtke et al.} 2023}).
Urbanization in particular, dramatically changes habitat structures and
increases the intensity of environmental stressors, influencing habitat
hydrology, temperature, and connectivity
(\citeproc{ref-hamer_amphibian_2008}{Hamer and McDonnell 2008};
\citeproc{ref-scheffers_effects_2012}{Scheffers and Paszkowski 2012};
\citeproc{ref-barrett_urbanization_2014}{Barrett and Price 2014};
\citeproc{ref-yang_global_2022}{Yang et al. 2022}). Amphibian residents
within urban habitats often display a reduction in body condition and
fitness when compared to rural populations
(\citeproc{ref-iglesias-carrasco_urban_2017}{Iglesias-Carrasco et al.
2017}; \citeproc{ref-jennette_variation_2019}{Jennette et al. 2019}).
Further, urban soil and water are prone to contamination from proximity
to industrial (\citeproc{ref-cachada_levels_2012}{Cachada et al. 2012};
\citeproc{ref-coban_soil_2022}{Coban et al. 2022};
\citeproc{ref-liu_microbial_2025}{Liu et al. 2025};
\citeproc{ref-mclellan_microbiome_2015}{McLellan et al. 2015}) and
agricultural pollution
(\citeproc{ref-soto-pozos_imprints_2025}{Soto-Pozos et al. 2025};
\citeproc{ref-jimenez_fungicide_2021}{Jiménez et al. 2021}),
subsequently serving as the vector through which amphibians are exposed
to pollutants (\citeproc{ref-dong_rice_2024}{Dong et al. 2024};
\citeproc{ref-purucker_estimating_2023}{Purucker et al. 2023}) and the
microbiota reservoir for dermal microbiota colonization
(\citeproc{ref-walke_amphibian_2014}{Walke et al. 2014};
\citeproc{ref-barnes_assembly_2021}{Barnes et al. 2021};
\citeproc{ref-loudon_microbial_2013}{Loudon et al. 2013},
\citeproc{ref-loudon_microbial_2013}{2013}).

To date, the majority of studies that have characterized amphibian
microbiota have done so through sequencing of the 16S rRNA marker gene
(\citeproc{ref-barnes_assembly_2021}{Barnes et al. 2021};
\citeproc{ref-santos_skin_2023}{Santos et al. 2023};
\citeproc{ref-muletz_wolz_effects_2018}{Muletz Wolz et al. 2018};
\citeproc{ref-afonso_effects_2022}{Afonso et al. 2022}). Metabarcoding
through amplicon enrichment utilizes existing microbial databases to
identify the taxonomy of prokaryotes in a manner that is generally
accessible and replicable (\citeproc{ref-pruesse_silva_2007}{Pruesse et
al. 2007}; \citeproc{ref-quast_silva_2013}{Quast et al. 2013}). While
suited for taxonomic identification of prokaryote taxa, amplicon
sequencing can result in a coarse taxonomic resolution and cannot
capture any of the functional variation that may occur within microbiome
communities (\citeproc{ref-campanaro_taxonomy_2018}{Campanaro et al.
2018}; \citeproc{ref-durazzi_comparison_2021}{Durazzi et al. 2021}).
While indirect methods to infer functional features of microbiomes from
16S sequences exist (\citeproc{ref-barberan_exploration_2012}{Barberán
et al. 2012}), these rely on genomes recorded in public databases, which
are largely biased towards humans and other model taxa. These databases
often have limited representation of amphibian-associated microbes,
which, despite similar 16S sequences, might carry significantly
different functional attributes
(\citeproc{ref-sun_inference-based_2020}{Sun et al. 2020};
\citeproc{ref-callahan_exact_2017}{Callahan et al. 2017}). By contrast,
genome-resolved metagenomics is a newer approach where random short-read
DNA sequences are used to assemble microbial genomes
(\citeproc{ref-weinstock_genomic_2012}{Weinstock 2012};
\citeproc{ref-quince_shotgun_2017}{Quince et al. 2017}), upon which
community-level microbiome inferences are built. This alternative
approach was considered preferable to 16S amplicon sequencing, as larger
sequence assemblies can improve taxonomic resolution
(\citeproc{ref-ziesemer_intrinsic_2015}{{Ziesemer et al.} 2015};
\citeproc{ref-durazzi_comparison_2021}{Durazzi et al. 2021}) and capture
functional gene information within coding regions of the microbial
genomes (\citeproc{ref-bates_microbiome_2022}{{Bates et al.} 2022};
\citeproc{ref-eisenhofer_impact_2023}{Eisenhofer et al. 2023}).
Additionally, because this methodology does not rely on primer-based
enrichment, which may result in biased representation-targeting
bacterial genome sequences, the relative abundance of taxa is better
preserved (\citeproc{ref-gloor_microbiome_2017}{Gloor et al. 2017}). The
compromise with this approach, however, is that host DNA will inevitably
be sequenced and will comprise the majority of the sequence reads
(\citeproc{ref-pereira-marques_impact_2019}{Pereira-Marques et al.
2019}). Removing host DNA generally entails mapping reads to a
reference, which in the case of amphibians are often unavailable or
difficult to construct owing to the often large size and complexity of
amphibian genomes (\citeproc{ref-sun_slow_2012}{Sun et al. 2012};
\citeproc{ref-liedtke_macroevolutionary_2018}{Liedtke et al. 2018};
\citeproc{ref-kosch_comparative_2025}{Kosch et al. 2025}).

One amphibian species whose dermal microbiome has received increased
attention, following outbreaks of \emph{Bsal} in Central Europe
(\citeproc{ref-castro_monzon_batrachochytrium_2022}{Castro Monzon et al.
2022}), are fire salamanders (\emph{Salamandra salamandra}; Linnaeus
1758). Widely dispersed and abundant through much of Europe, fire
salamanders are typically associated with forest biomes
(\citeproc{ref-wagner_habitat_2020}{Wagner et al. 2020};
\citeproc{ref-velo-anton_salamandra_2015}{Velo-Antón and Buckley 2015};
\citeproc{ref-dinis_dry-climate_2025}{Dinis et al. 2025}), but have also
been reported in a variety of urban environments throughout their
species range and can be considered tolerant to urban habitats
(\citeproc{ref-lourenco_trapped_2017}{Lourenço et al. 2017};
\citeproc{ref-iglesias-carrasco_urban_2017}{Iglesias-Carrasco et al.
2017}; \citeproc{ref-marques_coi_2022}{Marques et al. 2022};
\citeproc{ref-alarcon-rios_urban_2024}{Alarcón-Ríos et al. 2024}). Urban
habitats are generally associated with environmental perturbations to
microbial communities
(\citeproc{ref-mckinney_urbanization_2006}{McKinney 2006};
\citeproc{ref-zhou_urbanization_2023}{Zhou et al. 2023};
\citeproc{ref-zheng_urbanization_2025}{Zheng et al. 2025}). Previous
studies among urban amphibians report that urbanization corresponds to
increased homogeneity among urban populations but not necessarily
decreases in taxonomic richness or diversity
(\citeproc{ref-barnes_assembly_2021}{Barnes et al. 2021};
\citeproc{ref-basanta_genetic_2022}{Basanta et al. 2022};
\citeproc{ref-zhou_urbanization_2023}{Zhou et al. 2023};
\citeproc{ref-touceda-suarez_differences_2025}{Touceda-Suárez et al.
2025}; \citeproc{ref-wheeler_effects_2025}{Wheeler et al. 2025}) as
repeated disturbances may disrupt community connections providing
openings for colonization (\citeproc{ref-hu_impervious_2018}{Hu et al.
2018}).

Our aim in this study has been to characterize how dermal microbiomes
vary between habitats, by focusing on composition and functional
differences in the fire salamander dermal microbiome between urban and
rural habitats. To our knowledge, the fire salamander microbiome has
thus far only ever been subject to 16S amplicon sequencing, and while
characterization of the microbiome at variable life stages and habitats
are available (\citeproc{ref-sanchez_cutaneous_2017}{Sanchez et al.
2017}; \citeproc{ref-bletz_disruption_2018}{{Bletz et al.} 2018};
\citeproc{ref-afonso_effects_2022}{Afonso et al. 2022}), they miss out
on the more detailed taxonomic and functional resolution that shotgun
sequencing provides
(\citeproc{ref-antony-babu_multiple_2017}{Antony-Babu et al. 2017}).
Therefor, we applied shotgun-sequencing and assembled microbial genomes
to characterize the differences in bacterial composition and
functionality between fire salamanders from comparatively urban and
rural habitats. We hypothesized that the dermal microbiome of
salamanders would vary between urban and rural habitats in assemblage,
with the expectation that diversity and variance between sites would be
lower among urban sites owing to more recent anthropogenic disturbances
favouring generalist taxa tolerant to perturbations
(\citeproc{ref-zheng_urbanization_2025}{Zheng et al. 2025};
\citeproc{ref-clavel_worldwide_2011}{Clavel et al. 2011}). Furthermore,
that the inferred functional traits of microbiota assemblages would also
differ between habitats.

\section{Methods}\label{methods}

\subsubsection{Sampling}\label{sampling}

Field sampling was conducted at various urban parks in the city of Porto
in Portugal, Europe (Figure 1). Porto is a large urban matrix comprised
of isolated urban parks of varying size, extensive forested areas
comprised of either endemic (e.g.~\emph{Quercus, Acer}) and invasive
tree species (e.g.~\emph{Eucalyptus, Acacia}), and many small
agricultural plots intersecting them
(\citeproc{ref-guilherme_looking_2022}{Guilherme et al. 2022}). Sites
were generalized as either urban or rural based on their proximity to
the urban metropolitan core of the city
(\citeproc{ref-madureira_green_2011}{Madureira et al. 2011}), with up to
11 samples were collected when possible, with thirty samples collected
from each habitat. Urban sites were small semi-forested plots
(\textasciitilde9 ha) surrounded by areas of high impervious surface,
such as roads and built structures. Rural sites, by contrast, were large
areas (27-39 ha) with extensive tree cover, usually endemic deciduous
species, and surrounded by agricultural or stands of exotic tree
species. Individuals were collected between mid-October to early
November 2023, on nights of high humidity and when temperatures range
from 14-22 ºC. This coincides with a season of peak activity when males,
distinguished by the presence of an enlarged cloaca, move above the leaf
litter and present themselves as mating candidates
(\citeproc{ref-velo-anton_salamandra_2015}{Velo-Antón and Buckley
2015}). Using sterile gloves, specimens were collected and placed in
small individualized sterile containers, ideally avoiding direct
contact, instead scooping them into the containers to minimize contact
with the dermis. Salamanders were then moved to laboratory facilities at
Research Centre in Biodiversity and Genetic Resources (CIBIO) in Vairão,
Portugal, the following morning to proceed with swabbing. During
swabbing, a technician wearing sterile gloves (replaced between samples)
restrained individual salamanders by the base of the tail, using sterile
paper towel to help maintain a grip. Individual were then sprayed
thoroughly with distilled H\textsubscript{2}O to remove debris. Once
rinsed, the individual dermis was swabbed with ten stokes on the dorsal,
ventral, and each lateral side with a swab and placed in a 2 mL tube
pre-filled with DNA/RNA Shield (Zymo Research) and frozen at -20ºC.
After swabbing, individuals were returned to the sampled location later
the same day.

\subsubsection{Sequencing}\label{sequencing}

Samples were processed in October 2024 at the Globe Institute's
(University of Copenhagen) molecular laboratory in Copenhagen, Denmark,
following the established EHI laboratory protocols available at
www.earthhologenome.org/laboratory
(\citeproc{ref-pietroni_hologenomic_2025}{Pietroni et al. 2025}).
Samples underwent bead-beating before DNA isolation employing silica
magnetic beads with solid-phase reversible immobilization
(\citeproc{ref-lauritsen_robust_2025}{{Lauritsen et al.} 2025}). The
concentration of DNA extracts was quantified using a Qubit™ 3
Fluorometer (Thermo Fisher Scientific, USA). DNA was fragmented into
approximately 450 bp-long fragments using a Covaris LE220 platform.
Library preparation followed the ligation-based BEST protocol
(\citeproc{ref-caroe_single-tube_2018}{Carøe et al. 2018}), utilizing a
standard input of 200 ng of DNA in 24 µl or the closest amount feasible
based on the sample DNA concentration. We used 1.5 µl of 20 µM adapters
for a 50-200 ng DNA input, 1.5 µl of 10 µM for 10-50 ng, 1.5 µl of 5 µM
for \textless10 ng, and 1.5 µl of 2 µM for samples below the
quantification range of the Qubit® dsDNA HS (High Sensitivity) Assay
Kits. Libraries underwent qPCR screening to determine the optimal number
of library indexing PCR cycles
(\citeproc{ref-murray_benchtop_2015}{Murray et al. 2015}), followed by
PCR amplification using unique dual index primers with an adjusted
number of cycles ranging from X\textsubscript{1}-X\textsubscript{2}. The
resulting libraries underwent automated capillary electrophoresis
(Bioanalyzer or Fragment Analyzer) for assessment of fragment-length
distribution, adapter dimers, and adapter-to-library molar ratios.
Finally, samples were pooled into multiple sequencing batches, and
sequencing was performed across multiple lanes of NovaSeq6000 and
NovaSeq X platforms (Illumina, USA), generating 9 GB (approximately 30
million reads) of 150 bp paired-end sequencing data per sample.

\subsubsection{Bioinformatics}\label{bioinformatics}

Raw metagenomic reads were processed following the bioinformatic
workflow outlined by the Earth Hologenome Initiative (see {Leonard et
al.} (\citeproc{ref-leonard_global_2024}{2024});
\url{https://www.earthhologenome.org/}), and stored in the Electronic
Research Data Archive (ERDA) of the University of Copenhagen.
Computation was conducted in the local cluster of the Globe Institute
(Mjolnir) using the snakemake- and slurm-based bioinformatic pipelines
tailored to the platform
(\url{https://github.com/earthhologenome/EHI_bioinformatics/tree/mjolnir}).
Quality filtering to remove adapters, low-quality reads (phred
\textgreater= 20), short reads (\textgreater= 60 bp), and polyG/polyX
tails was done with Fastp v.0.23.4 (\citeproc{ref-chen_fastp_2018}{Chen
et al. 2018}). With skin swabs, the majority of DNA extracted is
expected to belong to the host, and so the proportion of prokaryotic
sequences were estimated using SingleM microbial fraction
(\citeproc{ref-eisenhofer_quantifying_2024}{Eisenhofer et al. 2024}).
Host-derived reads were removed by mapping against a reference host
genome using Bowtie2 v.2.5.1 (\citeproc{ref-langmead_fast_2012}{Langmead
and Salzberg 2012}) and Samtools v.1.18
(\citeproc{ref-danecek_twelve_2021}{{Danecek et al.} 2021}). In the
absence of a reference genome for \emph{S. salamandra}, we used a
\emph{Salamandra salamandra} transcriptome (NCBI RefSeq HBZU010000000;
Libro et al. (\citeproc{ref-libro_novo_2023}{2023})) in addition to the
reference genome of \emph{Pleurodeles waltl} (Michahelles, 1830) (NCBI
RefSeq assembly GCF\_031143425.1; {Brown et al.}
(\citeproc{ref-brown_chromosome-scale_2025}{2025})), which was the most
phylogenetically related species with a reference genome available at
the time (\citeproc{ref-zhang_phylogeny_2008}{Zhang et al. 2008}).
Unmapped metagenomic reads were assembled into contigs then into
putative bacterial genomes known as \emph{metagenome-assembled genomes
(MAGs)} by MEGAHIT v.1.2.9 (\citeproc{ref-li_megahit_2016}{Li et al.
2016}). The quality of these MAGs for subsequent downstream analyses was
evaluated with Quast (\citeproc{ref-gurevich_quast_2013}{Gurevich et al.
2013}) and contig coverage was determined by CoverM
(\citeproc{ref-aroney_coverm_2025}{Aroney et al. 2025}). Metagenomic
binning was performed by multiple algorithms, including CONCOCT v.1.1
(\citeproc{ref-alneberg_binning_2014}{Alneberg et al. 2014}), MetaBAT2
v.1.7 (\citeproc{ref-kang_metabat_2019}{Kang et al. 2019}), and MaxBin2
v.2.2.4 (\citeproc{ref-wu_maxbin_2016}{Wu et al. 2016}), the refined by
MetaWRAP (\citeproc{ref-uritskiy_metawrapflexible_2018}{Uritskiy et al.
2018}) into consensus bins. Bins that passed the CheckM2
(\citeproc{ref-chklovski_checkm2_2023}{Chklovski et al. 2023}) quality
criteria of \textgreater50\% completeness and \textless10\%
contamination were considered as metagenome-assembled genomes (MAGs) for
downstream analyses (\citeproc{ref-bowers_minimum_2017}{{Bowers et al.}
2017}).

Taxonomical annotations of MAGs were matched to the Genome Taxonomy
Database (GTDB) with the taxonomic classification toolkit (GTDB-Tk)
(\citeproc{ref-chaumeil_gtdb-tk_2022}{Chaumeil et al. 2022};
\citeproc{ref-parks_gtdb_2022}{Parks et al. 2022}). Genome-Inferred
Functional Traits (GIFTs) of MAGs were assigned with DRAM
(\citeproc{ref-shaffer_dram_2020}{{Shaffer et al.} 2020}) using the Pfam
(\citeproc{ref-finn_pfam_2016}{{Finn et al.} 2016}), KEGG
(\citeproc{ref-kanehisa_kegg_2017}{Kanehisa et al. 2017}), UniProt
(\citeproc{ref-mcgarvey_uniprot_2019}{McGarvey et al. 2019}), CAZY
(\citeproc{ref-cantarel_carbohydrate-active_2009}{Cantarel et al.
2009}), and MEROPS (\citeproc{ref-rawlings_merops_1999}{Rawlings and
Barrett 1999}) databases and distillR (available at
\url{https://github.com/anttonalberdi/distillR}). An average nucleotide
identity (ANI) of 98\% and a minimum genome completeness of 50\% were
used to dereplicate highly similar MAGs using dRep
(\citeproc{ref-olm_drep_2017}{Olm et al. 2017}). Reads were then mapped
to a MAG catalogue to extract quantitative read-abundance information to
calculate relative abundance of each MAG per sample.

\subsubsection{Statistics}\label{statistics}

For each sample, we determined the proportion of sequenced bases that
were mapped to MAGs as compared to those that were filtered due to poor
quality, unmapped, or mapped to the host. The domain-adjusted mapping
rate (DAMR) was used to quantify microbial genome recovery from
metagenomes, calculated as the rate of read mapping to sets of microbial
genomes (MAGs or other references) divided by the fraction of the
community predicted to be prokaryote
(\citeproc{ref-eisenhofer_quantifying_2024}{Eisenhofer et al. 2024}).
Differences in the proportion of mapped bases between habitats were
determined through a Wilcox rank test
(\citeproc{ref-bauer_constructing_1972}{Bauer 1972}). The relative
abundance of each MAG was defined as the number of assigned reads as a
proportion of the MAGs' length, divided by average read length
(\citeproc{ref-gloor_compositional_2016}{Gloor and Reid 2016};
\citeproc{ref-gloor_microbiome_2017}{Gloor et al. 2017}), with a centred
log ratio (CLR) transformation used to remove the constant-sum
constraint and to centre values around zero
(\citeproc{ref-aitchison_statistical_1982}{Aitchison 1982}). Differences
between habitats in the relative abundances of taxa were determined with
an Analysis of Compositions of Microbiomes with Bias Correction 2
(ANCOMB2) (\citeproc{ref-lin_analysis_2020}{Lin and Peddada 2020}) and
reported as log fold changes between habitats; this analysis was
repeated at each taxonomic resolution.

Differences in MAG richness as well as neutral, phylogenetic, and
functional beta diversities were calculated with Hill numbers
(\citeproc{ref-hill_diversity_1973}{Hill 1973};
\citeproc{ref-chao_unifying_2014}{Chao et al. 2014}) using Hilldiv2
v.2.0.2 (\citeproc{ref-alberdi_hilldiv2_2025}{Alberdi 2025}). Neutral
diversity represent species diversity without considering the degree of
relatedness among MAGs, phylogenetic diversity incorporate branch-length
information of the phylogenetic tree of the MAGs, and functional
diversity consider functional differences among MAGs
(\citeproc{ref-hernandez_contrasting_2024}{Hernández et al. 2024}).
Differences in richness and diversity estimates between habitats and
or/sites were determined with the \emph{betadisper} function in vegan
v.2.6-10 (\citeproc{ref-oksanen_package_2013}{{Oksanen et al.} 2013}). A
permutational multivariate analysis of variance (PERMANOVA) was
performed to test for differences in microbial composition between
habitats using the \emph{adonis2} function in vegan. Once again,
Wilcoxon rank sum test was used to test for differences in the variance
within each group when comparing habitats and an analysis of variance
(ANOVA; Chambers and Hastie
(\citeproc{ref-chambers_statistical_1992}{1992})) was used to compare
differences in variance between sites.

For each MAG, GIFTs were aggregated to a singular GIFT score for each
functional element, then aggregated to genome specific GIFTs for each
individual using the \emph{to.element} and \emph{to.community} functions
included in distillR. GIFT scores ranged from 0 to 1, representing
representing the proportion of genes required for the functional
pathway, where 0 represents an absence of associated genes and 1 a
complete presence of all genes
(\citeproc{ref-aizpurua_impact_2026}{Aizpurua et al. 2026}). Functional
elements with less than an average score of 0.1 were excluded from
consideration. Among the remainder, a Wilcoxon rank sum test was used to
determine if there were differences in GIFTs between habitats, with a
subsequent correction for the false discovery rate. Additionally, among
any microbial taxa which significantly differed in relative abundance
between habitats, we also tested for significant differences in
functional element association using the aforementioned methods. All
statistical analyses were performed using the R software v.4.5.0 (R Core
Team, 2025).

\section{Results}\label{results}

From each site, up to 11 samples were collected when possible, with
thirty samples collected from each habitat (Figure 1). Sequencing
libraries included a total of 60 swabs used to generate approximately
105M total reads (mean ± SD = 1.75 ± 2.69, median = 1.11, M reads per
sample) of short-read DNA sequencing data. Quality filtering and host
read removal resulted in 98.5M filtered reads (mean ± SD = 1.69 ± 2.68 ,
median = 0.991, Gb per sample) of metagenomic data, assembled into 118
metagenome-assembled genomes (MAGs) that passed our minimum quality
criteria (CheckM2 estimated completeness \textgreater50\%, contamination
\textless10\%). A mean of 4.42 ± 1.89 Gigabases were generated per
sample, with an average 5.8\% mapped to MAGs, while 5.3\% were removed
for low quality, and the remaining either unmapped or mapped to the
host. Average MAG completeness was 79.3 ± 16.1 \%, and contamination
2.75 ± 2.42 \%. The mean DAMR suggest that 22.5 ± 20.8 \% of bacterial
reads were successfully mapped to MAGs. Between the two habitats, a
significantly lower proportion of bases were assigned to MAGs among
urban samples as compared to rural samples (W = 646, p = 0.003).

\begin{figure}
\centering
\pandocbounded{\includegraphics[keepaspectratio,alt={Figure 1. Sites were designated as either Urban or Rural habitats, depending on the proportion in impervious surface nearby. The names and number of swabs collected from each site are listed at each point. Photo courtesy ©Max Benito and Jose Vieira from ExSitu Project.}]{ms_ssal_porto_microbiome_files/documentation/porto_sample_sites_count.png}}
\caption{Figure 1. Sites were designated as either Urban or Rural
habitats, depending on the proportion in impervious surface nearby. The
names and number of swabs collected from each site are listed at each
point. Photo courtesy ©Max Benito and Jose Vieira from ExSitu Project.}
\end{figure}

Among the 118 distinct MAGs, the three most taxa rich phyla were
Bacteroidota (33.1\%), Pseudomonadota (33.1\%), and Bacillota
(21.2\%)(Table 2; Figure 2). Except for Pseudomonadota, the majority of
MAGs (73.7\%) were identified as new species, with a total of 87 MAGs
lacking species level annotations and a further 19 MAGs lacked genus
level annotations. Among families with more than a single ascribed taxa,
two, Erysipelotrichaceae (phylum Bacillota; 3 MAGs) and
Pseudonocardiaceae (phylum Actinomycetota; 3 MAGs), were exclusive to
rural habitats, and another two, Chitinophagaceae (phylum Bacteroidota;
5 MAGs) and Steroidobacteraceae (phylum Pseudomonadota, 2 MAGs), were
exclusive to urban habitats.

Table 2. For each observed prokaryptic phylum, the number of distinct
MAGs assembled, the proportion of MAGs assigned to that phylum, and the
number of distinct classes, orders, and families in each, and the
proportion of MAGs identified as new species.

{\def\LTcaptype{none} % do not increment counter
\begin{longtable}[]{@{}
  >{\raggedright\arraybackslash}p{(\linewidth - 12\tabcolsep) * \real{0.2468}}
  >{\raggedleft\arraybackslash}p{(\linewidth - 12\tabcolsep) * \real{0.0779}}
  >{\raggedleft\arraybackslash}p{(\linewidth - 12\tabcolsep) * \real{0.1169}}
  >{\raggedleft\arraybackslash}p{(\linewidth - 12\tabcolsep) * \real{0.1169}}
  >{\raggedleft\arraybackslash}p{(\linewidth - 12\tabcolsep) * \real{0.1039}}
  >{\raggedleft\arraybackslash}p{(\linewidth - 12\tabcolsep) * \real{0.1299}}
  >{\raggedleft\arraybackslash}p{(\linewidth - 12\tabcolsep) * \real{0.2078}}@{}}
\toprule\noalign{}
\begin{minipage}[b]{\linewidth}\raggedright
phylum
\end{minipage} & \begin{minipage}[b]{\linewidth}\raggedleft
MAGs
\end{minipage} & \begin{minipage}[b]{\linewidth}\raggedleft
\% MAGs
\end{minipage} & \begin{minipage}[b]{\linewidth}\raggedleft
Classes
\end{minipage} & \begin{minipage}[b]{\linewidth}\raggedleft
Orders
\end{minipage} & \begin{minipage}[b]{\linewidth}\raggedleft
Families
\end{minipage} & \begin{minipage}[b]{\linewidth}\raggedleft
\% new species
\end{minipage} \\
\midrule\noalign{}
\endhead
\bottomrule\noalign{}
\endlastfoot
Acidobacteriota & 2 & 1.70 & 1 & 2 & 2 & 1.00 \\
Actinomycetota & 8 & 6.80 & 1 & 2 & 6 & 0.75 \\
Bacillota & 25 & 21.2 & 2 & 5 & 10 & 1.00 \\
Bacteroidota & 39 & 33.00 & 1 & 4 & 10 & 0.90 \\
Desulfobacterota & 1 & 0.85 & 1 & 1 & 1 & 1.00 \\
Pseudomonadota & 39 & 33.00 & 2 & 12 & 16 & 0.36 \\
Verrucomicrobiota & 4 & 3.40 & 2 & 2 & 2 & 1.00 \\
\end{longtable}
}

\newpage

\begin{figure}
\centering
\pandocbounded{\includegraphics[keepaspectratio,alt={Figure 2. For each branch in the MAG phylogeny, concentric rings indicate the corresponding phylum (see legend), whether the taxa were identified in an urban (purple) or rural (green) habitat, and whether the MAG likely represents a new species (grey).}]{ms_ssal_porto_microbiome_files/Ssal_Microbiota/figure/circular_tree_simplified_legend.png}}
\caption{Figure 2. For each branch in the MAG phylogeny, concentric
rings indicate the corresponding phylum (see legend), whether the taxa
were identified in an urban (purple) or rural (green) habitat, and
whether the MAG likely represents a new species (grey).}
\end{figure}

Pseudomonadota was by far the most abundant phylum represented followed
by Bacteroidota and Actinomycetota (Table 3). Among each phylum, only
Bacillota significantly differed in relative abundance between sites
being more abundant among rural sample. Further increasing taxonomic
resolution, we found this to remain the case until reaching the family
resolution, where upon three families of Pseudomonadota (Moraxellaceae,
Micrococcaceae, and an undefined family of order Burkholderiales) were
also found to be more abundant among rural samples. At the highest
taxonomic resolution we identified 12 MAGs that significantly differed
between habitats, all of which were more abundant among rural samples
(Figure 3B); nine of these corresponded to the genus
\emph{Acinetobacter}, two to \emph{Citrobacter}, and one from each
family of Mycoplasmoidaceae and the order Burkholderiales (Table SM1).

Table 3. For each phylum, the mean CLR of relative abundances of each
phylum are reported among urban, rural, and overall habitats.
Significant differences between habitats (p \textless{} 0.05) were only
observed among Bacillota.

\begin{table}
\centering
\begin{tblr}[         %% tabularray outer open
]                     %% tabularray outer close
{                     %% tabularray inner open
colspec={Q[]Q[]Q[]Q[]Q[]},
hline{2}={1-5}{solid, black, 0.05em},
hline{1}={1-5}{solid, black, 0.08em},
hline{10}={1-5}{solid, black, 0.08em},
}                     %% tabularray inner close
Phylum & Urban & Rural & Overall & p \\
Pseudomonadota & 5.94 & 6.030 & 5.990 & 0.411000 \\
Bacteroidota & 3.84 & 2.620 & 3.230 & 0.099400 \\
Actinomycetota & 2.54 & 2.970 & 2.750 & 0.274000 \\
Bacillota & -2.14 & 0.574 & -0.782 & 0.000002 \\
Bacillota_A & -1.45 & -2.210 & -1.830 & 0.467000 \\
Acidobacteriota & -2.88 & -2.480 & -2.680 & 0.285000 \\
Verrucomicrobiota & -3.34 & -3.330 & -3.330 & 0.912000 \\
Desulfobacterota & -3.17 & -3.520 & -3.350 & 0.636000 \\
\end{tblr}
\end{table}

\newpage

\begin{figure}
\centering
\pandocbounded{\includegraphics[keepaspectratio,alt={Figure 3. A) Differences in relative abundance after a centred log ratio (CLR) transformation between rural and urban habitats at each of the detected phylum. B) Mean relative abundances after a CLR transformation are plotted against LFD, where taxa with significant differences are coloured according to their phylum (yellow = Pseudomonadota, blue = Bacillota).}]{ms_ssal_porto_microbiome_files/Ssal_Microbiota/figure/relative_abundance.png}}
\caption{Figure 3. A) Differences in relative abundance after a centred
log ratio (CLR) transformation between rural and urban habitats at each
of the detected phylum. B) Mean relative abundances after a CLR
transformation are plotted against LFD, where taxa with significant
differences are coloured according to their phylum (yellow =
Pseudomonadota, blue = Bacillota).}
\end{figure}

No significant differences in either species richness, neutral
diversity, phylogentic diversity, or functional diversity estimates were
observed between either habitat (Figure 4A-D). An analysis of variance
did, however, identify significant differences between habitats and
sites at each of the aforementioned metrics (Figure 4E-H; Table SM2).
Further, within each of the habitats, significant differences in neutral
diversity estimates were observed among sites from either habitat
(Figure S1B,F) and in phylogenetic variation exclusively among urban
sites (Figure S1G).

\newpage

\begin{figure}
\centering
\pandocbounded{\includegraphics[keepaspectratio,alt={Figure 4. Differences in (A-D) alpha within and (E-H) beta diversity between habitats when considering (A,F) species richness, (B,F) neutral diversity, (C,G) phylogenetic diversity, and (D,H) functional diversity.}]{ms_ssal_porto_microbiome_files/Ssal_Microbiota/figure/habitat_alphabeta_div.png}}
\caption{Figure 4. Differences in (A-D) alpha within and (E-H) beta
diversity between habitats when considering (A,F) species richness,
(B,F) neutral diversity, (C,G) phylogenetic diversity, and (D,H)
functional diversity.}
\end{figure}

After aggregating the functional association of microbial assemblages
among individuals from either habitat, the most complete functional
traits were those associated with biosynthesis of various organic
anions, amino, and nucleic acids (Figure S2). Significant differences in
the aggregare GIFT scores between habitat related to the biosynthesis of
citrate, degradation of tryptophan, and structural formation of
lipopolysaccharide between habitats (Figure 5). In each case, predicted
GIFT associations were greater among rural samples when compared to
urban counterparts. When comparing function element associations between
those microbial taxa that significantly differed between habitats and
those that did not, we observed significant differences in function
associations toward the biosynthesis of betaine (W = 955, p = 0.02), the
degradation of triglyceride lipids (W = 328, p = 0.02), lactose sugars
(W = 959, p = 0.02), and xenobiotic toluene (W = 264, p = 0.02), along
with differences in cellular structure associated with peptidoglycan (W
= 360, p = 0.04) and flagellum (W = 359, p = 0.04).

\newpage

\begin{figure}
\centering
\pandocbounded{\includegraphics[keepaspectratio,alt={Figure 5. Mean GIFT probability of functional elements that significantly differed between habitats.}]{ms_ssal_porto_microbiome_files/Ssal_Microbiota/figure/gift_class_habitat_comparisons.png}}
\caption{Figure 5. Mean GIFT probability of functional elements that
significantly differed between habitats.}
\end{figure}

\section{Discussion}\label{discussion}

While our hypothesis, that urbanization would result in lower taxonomic
and functional variance among sites, was not supported, significant
differences in diversity between sites and habitats were observed,
suggesting a shift in microbiome assemblage between habitats and an
increased phylogenetic diversity among urban salamanders microbiota.
Comparing the relative abundances of microbiota between habitats,
several taxa, namely from the phyla Pseudomonadota and Bacillia, were
significantly less abundant among urban salamanders. Furthermore, when
comparing significant differences in the inferred functional elements of
microbiome assemblages between habitats, urban microbiota assemblages
were found to be less likely associated with several functions when
compared to rural assemblages. Based on our results, it appears that
fire salamander residency in urbanized habitats coincides with a decline
in the relative abundance of several relatively common taxa, which may
consequentially result in decrease in functional elements possibly
associated with those taxa.

\subsection{Characterization of dermal
microbiota}\label{characterization-of-dermal-microbiota}

Comparing the application of shotgun sequencing in this study to
previously reported results from other studies using 16S amplicon
sequencing of the fire salamander dermal microbiome, both suggest a
similar microbiome assemblage where Pseudomonadota has been identified
as among the most taxa rich and abundant phylum, followed by
Bacteroidota and Actinomycetota
(\citeproc{ref-afonso_effects_2022}{Afonso et al. 2022};
\citeproc{ref-barnes_assembly_2021}{Barnes et al. 2021};
\citeproc{ref-bletz_host_2017}{Bletz et al. 2017},
\citeproc{ref-bletz_disruption_2018}{2018}). Unlike many of those
studies, however, Bacillota and Bacteroidota were identified as
similarly taxonomic rich phylum. This may reflect that unlike
Pseudomonadota, that MAGs ascribed to this phylum have comparatively few
existing taxonomic references. While studies utilizing 16S amplicon
sequences also report relative abundance, this measure is generated
through statistical simulation or inferences to correct for enrichment
bias (\citeproc{ref-mandal_analysis_2015}{Mandal et al. 2015};
\citeproc{ref-kurtz_sparse_2015}{Kurtz et al. 2015}). Amphibians
regularly shed their epidermis, in a process called sloughing
(\citeproc{ref-ohmer_skin_2015}{Ohmer et al. 2015}) which maintains a
healthy turnover of skin but, as revealed by electron microscopy and
qPCR (\citeproc{ref-bletz_disruption_2018}{{Bletz et al.} 2018}),
results in low microbial abundances along the skins surface. With less
starting genetic material, the number of enrichment cycles required for
suitable detection 16S fragments increases (30-35 cycles) as does the
subsequent risk of initial enrichment bias
(\citeproc{ref-nearing_identifying_2021}{Nearing et al. 2021};
\citeproc{ref-silverman_measuring_2019}{Silverman et al. 2019}). Shotgun
sequencing, by not relying on abundant enrichment cycles (3-10 cycles),
therefore does not need to infer relative abundance and is less
susceptible to starting material biases. The greatest logistical
challenge of this approach, particularly when applied to amphibians, was
the presence of a large complex host genome which required prior
filtering to separate host from bacterial DNA
(\citeproc{ref-eisenhofer_quantifying_2024}{Eisenhofer et al. 2024}). We
were able to assign an average of 5.8\% of reads per sample to bacterial
genomes and despite excluding the majority of reads, were able to
construct 118 bacterial genomes, the majority of which lack prior
taxonomic references. With an average genome size of 3.9Mb corresponding
to estimated 78.5\% completeness, we wish to highlight that with even a
small fraction of sequence reads, it is possible to generate
comprehensive bacterial genome assemblies. Additionally, while not
included in this study, the majority of reads attributed to the fire
salamander genome are available for subsequent analysis of host genomics
and will be a valuable contribution to amphibian genomic development
(\citeproc{ref-pietroni_hologenomic_2025}{Pietroni et al. 2025};
\citeproc{ref-kosch_amphibian_2024}{{Kosch et al.} 2024}).

A further benefit of implementing a shotgun sequencing approach was our
ability to predict microbiome functionality. Among similarly related
taxa, functionality can vary and will not be captured by 16S sequencing
unless the taxa in question have already been functionally characterized
(\citeproc{ref-cevallos_genomic_2022}{Cevallos et al. 2022}). Prior
studies comparing 16S community composition to functional gene
characterizations have found inconsistent associations between the two
(\citeproc{ref-barberan_exploration_2012}{Barberán et al. 2012}).
Shotgun sequencing, by capturing full genes, can then use GIFTs rather
than by taxonomic inference, which would be especially fraught given the
high number of taxa described in this study without prior taxonomic
references. Much of the work on microbiome functionality among
amphibians, and in particular fire salamanders, have been to determine
the role of microbiota in conferring resistance to chytrid fugal
implications (\citeproc{ref-bletz_disruption_2018}{{Bletz et al.} 2018};
\citeproc{ref-bates_captivity_2019}{Bates et al. 2019};
\citeproc{ref-crawshaw_widespread_2022}{Crawshaw et al. 2022};
\citeproc{ref-smith_skin_2018}{Smith et al. 2018};
\citeproc{ref-kueneman_effects_2022}{{Kueneman et al.} 2022}). Among the
microbiome communities, the genera \emph{Pseudomonas} (mean CLR = 72)
followed by \emph{Acinetobacter} (mean CLR = 34) were the most abundant
among all samples. The relative abundance of \emph{Pseudomonas,}
however, did not vary significantly between habitats. Rather differences
in microbiome abundances between habitats were driven by the relative
abundance of specific taxa, most notably the genera
\emph{Acinetobacter}, \emph{Citrobacter}, and \emph{Mycoplasma}. Among
these, \emph{Acinetobacter} has repeatedly been identified for it's
associated with inhibiting \emph{Bsal} and \emph{Bd} growth
(\citeproc{ref-muletz-wolz_inhibition_2017}{Muletz-Wolz et al. 2017};
\citeproc{ref-cevallos_genomic_2022}{Cevallos et al. 2022};
\citeproc{ref-hoyt_bacteria_2015}{Hoyt et al. 2015}). Meanwhile,
\emph{Citrobacter} is mostly known as a foodborne pathogen associated
with aquaculture (\citeproc{ref-yang_comparative_2024}{Yang et al.
2024}; \citeproc{ref-paranhos_sepsis_2024}{Paranhos et al. 2024};
\citeproc{ref-li_damage_2022}{Li et al. 2022}) and \emph{Mycoplasma} is
generally recognized as a pathogen among reptiles, associated with
respiratory illness (\citeproc{ref-faulhaber_detection_2025}{Faulhaber
et al. 2025}; \citeproc{ref-brown_mycoplasma_2018}{Brown et al. 2018}).

Among the inferred functional traits identified as differing between
habitats we identified the biosynthesis of citrate, which is often
produced in the mitochondria and considered a product of metabolic
activity (\citeproc{ref-iacobazzi_citrate_2014}{Iacobazzi and Infantino
2014}); the degradation of tryptophan, an amino-acid involved in the
biosynthesis of proteins and usually converted into various metabolites
(\citeproc{ref-friedman_analysis_2018}{Friedman 2018}); and the
structural development of lipopolysaccharide, a component of
gram-negative bacterial cell walls that are generally considered
pathogenic (\citeproc{ref-mayeux_pathobiology_1997}{Mayeux 1997};
\citeproc{ref-raetz_lipopolysaccharide_2002}{Raetz and Whitfield 2002})
and in some cases used to explore immune responses among exposed
amphibians (\citeproc{ref-mcnabb_effects_2019}{McNabb et al. 2019};
\citeproc{ref-ferreira_lipopolysaccharide_2021}{Ferreira et al. 2021}).
Regrettably, while we identified these functional traits, literature on
their relevance to amphibian dermal health is limited, and deciphering
the role would require experimental study beyond the scope of this
study. By identifying functionally variability of the microbiome, future
inquiry can identify traits to examine the role that microbiota have in
amphibians' dermal health and pathology.

\subsection{Dermal microbiome differences between
habitats}\label{dermal-microbiome-differences-between-habitats}

Similar to other characterizations of amphibian dermal microbiota in
other studies (\citeproc{ref-barnes_assembly_2021}{Barnes et al. 2021};
\citeproc{ref-walke_seasonal_2021}{Walke et al. 2021}), diversity among
fire salamanders was not necessarily reduced among urban habitats, but
instead varied in composition. All observed phyla were detected in
either habitat, and among the few cases where a taxonomic family was
exclusive to a habitat, as was the case with Chitinophagaceae among
urban habitats and Pseudonocardiaceae and Erysipelotrichaceae among
rural habitats, they were nominally represented in regard to relative
abundance, contributing no more than 2\% of the total reads. Instead, it
appears that urban samples had a greater diversity of lineages, given
the increased phylogenetic diversity observed between urban sites. This
is in line with studies of urban soil microbiomes, which reported
greater phylogenetic diversity among urban sites
(\citeproc{ref-touceda-suarez_differences_2025}{Touceda-Suárez et al.
2025}). Specifically, a network analysis of urban soil microbiota in
Chengdu, China by Zhou et al.
(\citeproc{ref-zhou_urbanization_2023}{2023}) reported that soil in
urbanized sites were more phylogenetically diverse but subsequently had
lower taxonomic co-occurrence compared to rural sites. Changes in
microbial assemblages of the environment may occur through translocation
or deposition of soils between urban sites or through pollution
(\citeproc{ref-cachada_levels_2012}{Cachada et al. 2012};
\citeproc{ref-touceda-suarez_differences_2025}{Touceda-Suárez et al.
2025}). While amphibians' dermal microbiomes are generally distinct from
their environments (\citeproc{ref-burgon_phylogenomic_2021}{{Burgon et
al.} 2021}; \citeproc{ref-bletz_host_2017}{Bletz et al. 2017}), their
environment represents a reservoir for exposure to microbiota taxa
(\citeproc{ref-muletz_soil_2012}{Muletz et al. 2012}). With greater
phylogenetic variation between sites, more evenly distributed relative
abundance of taxa, and lower sequence output from urban salamander
microbiomes, the characterization of dermal microbiota in urban
salamander appears consistent with a microbial community that reduced
taxonomic redundancy and susceptible to stochastic disturbance.

\subsection{Conclusion}\label{conclusion}

Through the application of genome-resolved metagenomics, we were able to
expand upon the existing characterization of the fire salamander dermal
microbiome. We found support for variance between urban and rural
habitats in assemblage; however, this applied primarily to composition
of community assemblage and not microbiota diversity therein. The most
abundant taxa identified were consistent with relative abundances
estimated from previous studies, but many rare or low-abundance taxa
were also identified, which may be putatively new species. Furthermore,
urbanization has an apparent influence on the relative abundance of
several microbial taxa communities, specifically those that have
reported pathogenic or antifungal properties. Specifically, the
increased relative abundance of the presumed fungal-resistant
\emph{Acinetobacter} along with other taxa associated with pathogenic
impacts, may suggest greater exposure among rural salamanders to
potentially harmful microbial elements. Further study is recommended to
compare these impacts that altered bacterial colonies have on dermal
health and whether these changes better reflect variation in the
microbial composition of the habitat or a response to a yet identified
environmental pressure.

\newpage

\section{Acknowledgements}\label{acknowledgements}

Thanks to Carlos Caballero-Díaz for assistance with sampling, to Isolde
van Riemsdijk for consulting the best practices for sampling dermal
swabs.

This publication was funded through the generous support of the Fundação
para a Ciência e a Tecnologia (FCT) as part of the project
``ANTHROPHIBIANS: adaptive genomic and microbiotic responses of
amphibians in anthropogenic landscapes'' (2021.05570.BD;
\url{https://doi.org/10.54499/2021.05570.BD}).

\subsubsection{Data accessibility}\label{data-accessibility}

The raw data tables containing the quantitative information and
bioinformatic code scripts used for the analysis are available in a
dedicated GitHub repository accessible through Zenodo with Doi:
10.5281/zenodo.19578779. Raw sequencing data and microbial genome
sequences, have been or will be published as part of the 2nd EHI Data
Release, under the Bioproject PRJEB77033.

\newpage

\section{Supplementary Materials}\label{supplementary-materials}

\begin{figure}
\centering
\pandocbounded{\includegraphics[keepaspectratio,alt={Figure S1. Differences in beta diversity of species richness (A,E), Shannon diversity (C,F), phylogentic diversity (C,G), and functional diversity (D,H) between sites within rural (A-D) and urban (E-H) habitats.}]{ms_ssal_porto_microbiome_files/Ssal_Microbiota/figure/site_beta_diversity.png}}
\caption{Figure S1. Differences in beta diversity of species richness
(A,E), Shannon diversity (C,F), phylogentic diversity (C,G), and
functional diversity (D,H) between sites within rural (A-D) and urban
(E-H) habitats.}
\end{figure}

\begin{figure}
\centering
\pandocbounded{\includegraphics[keepaspectratio,alt={Figure S2. For each samples, the probability that the dermal microbiome, as an aggregate, is associated with a genome-inferred functional trait (GIFT) associated with biosynthesis, degradation, or structural (Str.) functionality.}]{ms_ssal_porto_microbiome_files/Ssal_Microbiota/figure/habitat_gift_function.png}}
\caption{Figure S2. For each samples, the probability that the dermal
microbiome, as an aggregate, is associated with a genome-inferred
functional trait (GIFT) associated with biosynthesis, degradation, or
structural (Str.) functionality.}
\end{figure}

Table S1. A list of MAGs for which a significant LFC was observed
between habitats, with negatives values indicating greater relative
abundance among rural samples.

\begin{table}
\centering
\begin{tblr}[         %% tabularray outer open
]                     %% tabularray outer close
{                     %% tabularray inner open
colspec={Q[]Q[]Q[]Q[]Q[]Q[]Q[]Q[]Q[]},
hline{2}={1-9}{solid, black, 0.05em},
hline{1}={1-9}{solid, black, 0.08em},
hline{14}={1-9}{solid, black, 0.08em},
}                     %% tabularray inner close
Phylum & Class & Order & Family & Genus & Species & MAG & LFC & p \\
Bacillota & Bacilli & Mycoplasmatales & Mycoplasmoidaceae & Mycoplasma_N &  & EHA05315_bin.4 & -4.5 & 9.4e-06 \\
Pseudomonadota & Gammaproteobacteria & Burkholderiales & 2013Ark19i &  &  & EHA05315_bin.2 & -3.8 & 2.7e-04 \\
Pseudomonadota & Gammaproteobacteria & Pseudomonadales & Moraxellaceae & Acinetobacter & Acinetobacter ursingii & EHA05315_bin.14 & -3.7 & 8.4e-06 \\
Pseudomonadota & Gammaproteobacteria & Pseudomonadales & Moraxellaceae & Acinetobacter & Acinetobacter johnsonii & EHA05315_bin.31 & -3.2 & 9.0e-05 \\
Pseudomonadota & Gammaproteobacteria & Pseudomonadales & Moraxellaceae & Acinetobacter & Acinetobacter bereziniae & EHA05315_bin.118 & -3.1 & 8.2e-05 \\
Pseudomonadota & Gammaproteobacteria & Pseudomonadales & Moraxellaceae & Acinetobacter & Acinetobacter sp001696615 & EHA05315_bin.26 & -2.5 & 1.5e-03 \\
Pseudomonadota & Gammaproteobacteria & Pseudomonadales & Moraxellaceae & Acinetobacter & Acinetobacter pittii_E & EHA05315_bin.3 & -2.5 & 1.1e-03 \\
Pseudomonadota & Gammaproteobacteria & Pseudomonadales & Moraxellaceae & Acinetobacter & Acinetobacter vivianii & EHA05315_bin.7 & -2.4 & 2.3e-03 \\
Pseudomonadota & Gammaproteobacteria & Pseudomonadales & Moraxellaceae & Acinetobacter & Acinetobacter beijerinckii & EHA05315_bin.52 & -2.4 & 1.2e-03 \\
Pseudomonadota & Gammaproteobacteria & Pseudomonadales & Moraxellaceae & Acinetobacter & Acinetobacter sp002367455 & EHA05315_bin.29 & -2.4 & 4.9e-03 \\
Pseudomonadota & Gammaproteobacteria & Enterobacterales & Enterobacteriaceae & Citrobacter & Citrobacter braakii & EHA05315_bin.84 & -2.1 & 4.4e-03 \\
Pseudomonadota & Gammaproteobacteria & Enterobacterales & Enterobacteriaceae & Citrobacter & Citrobacter freundii & EHA05315_bin.78 & -1.8 & 1.2e-02 \\
\end{tblr}
\end{table}

Table S1. ANOVA results comparing metrics for species richness, neutral
diversity, phylogenetic diversity, and functional diversity between
habitats. R\textsuperscript{2} represents the proportional sum of
squares and the significance.

\begin{table}
\centering
\begin{tblr}[         %% tabularray outer open
]                     %% tabularray outer close
{                     %% tabularray inner open
colspec={Q[]Q[]Q[]Q[]Q[]},
hline{2}={1-5}{solid, black, 0.05em},
hline{1}={1-5}{solid, black, 0.08em},
hline{6}={1-5}{solid, black, 0.08em},
}                     %% tabularray inner close
Metric & Habitat R2 & Habitat P & Site R2 & Site P \\
species richness & 0.03870089 & 0.0370 & 0.1280550 & 0.0336 \\
neutral diversity & 0.08861909 & 0.0003 & 0.2326501 & 0.0001 \\
phylogenetic diversity & 0.09280180 & 0.0001 & 0.2858097 & 0.0001 \\
functional diversity & 0.11158579 & 0.0028 & 0.2003820 & 0.0113 \\
\end{tblr}
\end{table}

\newpage

\section*{Bibliography}\label{bibliography}
\addcontentsline{toc}{section}{Bibliography}

\protect\phantomsection\label{refs}
\begin{CSLReferences}{1}{1}
\bibitem[\citeproctext]{ref-afonso_effects_2022}
Afonso, Mariana, Laura Coelho, Fátima Jesus, et al. 2022. {``Effects of
{Pine} and {Eucalypt} Ashes on Bacterial Isolates from the Skin
Microbiome of the Fire Salamander ({Salamandra} Salamandra).''}
\emph{Science of The Total Environment} 841: 156677.
\url{https://doi.org/10.1016/j.scitotenv.2022.156677}.

\bibitem[\citeproctext]{ref-aitchison_statistical_1982}
Aitchison, J. 1982. {``The {Statistical} {Analysis} of {Compositional}
{Data}.''} \emph{Journal of the Royal Statistical Society: Series B
(Methodological)} 44 (2): 139--60.
\url{https://doi.org/10.1111/j.2517-6161.1982.tb01195.x}.

\bibitem[\citeproctext]{ref-aizpurua_impact_2026}
Aizpurua, Ostaizka, Evie Morris, Lasse Nyholm, et al. 2026. {``Impact of
Sequencing Approach on Ecological Inference from Highly Variable Bat Gut
Microbiomes.''} \emph{bioRxiv}, January, 2026.01.08.698297.
\url{https://doi.org/10.64898/2026.01.08.698297}.

\bibitem[\citeproctext]{ref-akat_comden_complex_2023}
Akat Çömden, Esra, Melodi Yenmiş, and Berna Çakır. 2023. {``The Complex
Bridge Between Aquatic and Terrestrial Life: {Skin} Changes During
Development of Amphibians.''} \emph{Journal of Developmental Biology} 11
(1): 6. \url{https://doi.org/10.3390/jdb11010006}.

\bibitem[\citeproctext]{ref-alarcon-rios_urban_2024}
Alarcón-Ríos, Lucía, Antigoni Kaliontzopoulou, David Álvarez, and
Guillermo Velo-Antón. 2024. {``Urban Life Affects Differentiation and
Phenotypic Variation but Not Asymmetry in a Fully Terrestrial
Salamander.''} \emph{Evolutionary Biology} 51 (3): 295--309.
\url{https://doi.org/10.1007/s11692-024-09635-6}.

\bibitem[\citeproctext]{ref-alberdi_hilldiv2_2025}
Alberdi, Antton. 2025. \emph{Hilldiv2: {Integral} {Analysis} of
{Diversity} {Based} on {Hill} {Numbers}}.
\url{https://github.com/anttonalberdi/hilldiv2}.

\bibitem[\citeproctext]{ref-alneberg_binning_2014}
Alneberg, Johannes, Brynjar Smári Bjarnason, Ino De Bruijn, et al. 2014.
{``Binning Metagenomic Contigs by Coverage and Composition.''}
\emph{Nature Methods} 11: 1144--46.
\url{https://doi.org/10.1038/nmeth.3103}.

\bibitem[\citeproctext]{ref-antony-babu_multiple_2017}
Antony-Babu, Sanjay, Didier Stien, Véronique Eparvier, Delphine Parrot,
Sophie Tomasi, and Marcelino T Suzuki. 2017. {``Multiple {Streptomyces}
Species with Distinct Secondary Metabolomes Have Identical {16S} {rRNA}
Gene Sequences.''} \emph{Scientific Reports} 7: 11089.
\url{https://doi.org/10.1038/s41598-017-11363-1}.

\bibitem[\citeproctext]{ref-aroney_coverm_2025}
Aroney, Samuel TN, Rhys JP Newell, Jakob N Nissen, Antonio Pedro
Camargo, Gene W Tyson, and Ben J Woodcroft. 2025. {``{CoverM}: {Read}
Alignment Statistics for Metagenomics.''} \emph{Bioinformatics} 41 (4):
btaf147. \url{https://doi.org/10.1093/bioinformatics/btaf147}.

\bibitem[\citeproctext]{ref-barberan_exploration_2012}
Barberán, Albert, Antoni Fernández-Guerra, Brendan JM Bohannan, and
Emilio O Casamayor. 2012. {``Exploration of Community Traits as
Ecological Markers in Microbial Metagenomes.''} \emph{Molecular Ecology}
21 (8): 1909--17.
\url{https://doi.org/10.1111/j.1365-294X.2011.05383.x}.

\bibitem[\citeproctext]{ref-barnes_assembly_2021}
Barnes, Elle M, Steve Kutos, Nina Naghshineh, Marissa Mesko, Qing You,
and JD Lewis. 2021. {``Assembly of the Amphibian Microbiome Is
Influenced by the Effects of Land-Use Change on Environmental
Reservoirs.''} \emph{Environmental Microbiology} 23 (8): 4595--611.
\url{https://doi.org/10.1111/1462-2920.15653}.

\bibitem[\citeproctext]{ref-barrett_urbanization_2014}
Barrett, Kyle, and Steven J Price. 2014. {``Urbanization and Stream
Salamanders: A Review, Conservation Options, and Research Needs.''}
\emph{Freshwater Science} 33 (3): 927--40.
\url{https://doi.org/10.1086/677556}.

\bibitem[\citeproctext]{ref-basanta_genetic_2022}
Basanta, María Delia, Eria A Rebollar, Mirna G García-Castillo, et al.
2022. {``Genetic Variation of {Batrachochytrium} Dendrobatidis Is Linked
to Skin Bacterial Diversity in the {Pacific} Treefrog {Hyliola} Regilla
(Hypochondriaca).''} \emph{Environmental Microbiology} 24 (1): 494--506.
\url{https://doi.org/10.1111/1462-2920.15861}.

\bibitem[\citeproctext]{ref-bates_captivity_2019}
Bates, Kieran A, Jennifer MG Shelton, Victoria L Mercier, et al. 2019.
{``Captivity and Infection by the Fungal Pathogen {Batrachochytrium}
Salamandrivorans Perturb the Amphibian Skin Microbiome.''}
\emph{Frontiers in Microbiology} 10: 1834.
\url{https://doi.org/10.3389/fmicb.2019.01834}.

\bibitem[\citeproctext]{ref-bates_microbiome_2022}
{Bates, Kieran A, Ulf Sommer, Kevin P Hopkins, et al.} 2022.
{``Microbiome Function Predicts Amphibian Chytridiomycosis Disease
Dynamics.''} \emph{Microbiome} 10 (1): 44.
\url{https://doi.org/10.1186/s40168-021-01215-6}.

\bibitem[\citeproctext]{ref-bauer_constructing_1972}
Bauer, David F. 1972. {``Constructing Confidence Sets Using Rank
Statistics.''} \emph{Journal of the American Statistical Association} 67
(339): 687--90. \url{https://doi.org/10.1080/01621459.1972.10481279}.

\bibitem[\citeproctext]{ref-becker_cutaneous_2010}
Becker, Matthew H., and Reid N. Harris. 2010. {``Cutaneous {Bacteria} of
the {Redback} {Salamander} {Prevent} {Morbidity} {Associated} with a
{Lethal} {Disease}.''} \emph{PLoS One} 5 (6): 1--6.
\url{https://doi.org/10.1371/journal.pone.0010957}.

\bibitem[\citeproctext]{ref-bletz_host_2017}
Bletz, Molly C, Holly Archer, Reid N Harris, et al. 2017. {``Host
Ecology Rather Than Host Phylogeny Drives Amphibian Skin Microbial
Community Structure in the Biodiversity Hotspot of {Madagascar}.''}
\emph{Frontiers in Microbiology} 8: 1530.
\url{https://doi.org/10.3389/fmicb.2017.01530}.

\bibitem[\citeproctext]{ref-bletz_disruption_2018}
{Bletz, Molly C, Moira Kelly, Joana Sabino-Pinto, et al.} 2018.
{``Disruption of Skin Microbiota Contributes to Salamander Disease.''}
\emph{Proceedings of the Royal Society B} 285 (1885): 20180758.
\url{https://doi.org/10.1098/rspb.2018.0758}.

\bibitem[\citeproctext]{ref-bowers_minimum_2017}
{Bowers, Robert M, Nikos C Kyrpides, Ramunas Stepanauskas, et al.} 2017.
{``Minimum Information about a Single Amplified Genome ({MISAG}) and a
Metagenome-Assembled Genome ({MIMAG}) of Bacteria and Archaea.''}
\emph{Nature Biotechnology} 35 (8): 725--31.
\url{https://doi.org/10.1038/nbt.3893}.

\bibitem[\citeproctext]{ref-brown_mycoplasma_2018}
Brown, Daniel R., Meghan May, Janet M. Bradbury, et al. 2018.
{``Mycoplasma.''} In \emph{Bergey's {Manual} of {Systematics} of
{Archaea} and {Bacteria}}.
\url{https://doi.org/10.1002/9781118960608.gbm01263.pub2}.

\bibitem[\citeproctext]{ref-brown_chromosome-scale_2025}
{Brown, Thomas, Ketan Mishra, Ahmed Elewa, et al.} 2025.
{``Chromosome-Scale Genome Assembly Reveals How Repeat Elements Shape
Non-Coding {RNA} Landscapes Active During Newt Limb Regeneration.''}
\emph{Cell Genomics} 5 (2).
\url{https://doi.org/10.1016/j.xgen.2025.100761}.

\bibitem[\citeproctext]{ref-burgon_phylogenomic_2021}
{Burgon, James D, Miguel Vences, Sebastian Steinfartz, et al.} 2021.
{``Phylogenomic Inference of Species and Subspecies Diversity in the
{Palearctic} Salamander Genus {Salamandra}.''} \emph{Molecular
Phylogenetics and Evolution} 157: 107063.
\url{https://doi.org/10.1016/j.ympev.2020.107063}.

\bibitem[\citeproctext]{ref-cachada_levels_2012}
Cachada, A, P Pato, T Rocha-Santos, E Ferreira da Silva, and AC Duarte.
2012. {``Levels, Sources and Potential Human Health Risks of Organic
Pollutants in Urban Soils.''} \emph{Science of the Total Environment}
430: 184--92. \url{https://doi.org/10.1016/j.scitotenv.2012.04.075}.

\bibitem[\citeproctext]{ref-callahan_exact_2017}
Callahan, Benjamin J, Paul J McMurdie, and Susan P Holmes. 2017.
{``Exact Sequence Variants Should Replace Operational Taxonomic Units in
Marker-Gene Data Analysis.''} \emph{The ISME Journal} 11 (12): 2639--43.
\url{https://doi.org/10.1038/ismej.2017.119}.

\bibitem[\citeproctext]{ref-campanaro_taxonomy_2018}
Campanaro, Stefano, Laura Treu, Panagiotis G Kougias, Xinyu Zhu, and
Irini Angelidaki. 2018. {``Taxonomy of Anaerobic Digestion Microbiome
Reveals Biases Associated with the Applied High Throughput Sequencing
Strategies.''} \emph{Scientific Reports} 8 (1): 1926.
\url{https://doi.org/10.1038/s41598-018-20414-0}.

\bibitem[\citeproctext]{ref-cantarel_carbohydrate-active_2009}
Cantarel, Brandi L, Pedro M Coutinho, Corinne Rancurel, Thomas Bernard,
Vincent Lombard, and Bernard Henrissat. 2009. {``The
{Carbohydrate}-{Active} {EnZymes} Database ({CAZy}): An Expert Resource
for Glycogenomics.''} \emph{Nucleic Acids Research} 37 (suppl\_1):
D233--38. \url{https://doi.org/10.1093/nar/gkn663}.

\bibitem[\citeproctext]{ref-caroe_single-tube_2018}
Carøe, Christian, Shyam Gopalakrishnan, Lasse Vinner, et al. 2018.
{``Single-Tube Library Preparation for Degraded {DNA}.''} \emph{Methods
in Ecology and Evolution} 9 (2): 410--19.
\url{https://doi.org/10.1111/2041-210X.12871}.

\bibitem[\citeproctext]{ref-castro_monzon_batrachochytrium_2022}
Castro Monzon, Federico, Mark-Oliver Rödel, Florian Ruland, Gabriela
Parra-Olea, and Jonathan M Jeschke. 2022. {``Batrachochytrium
Salamandrivorans' Amphibian Host Species and Invasion Range.''}
\emph{EcoHealth} 19 (4): 475--86.
\url{https://doi.org/10.1007/s10393-022-01620-9}.

\bibitem[\citeproctext]{ref-cevallos_genomic_2022}
Cevallos, MA, MD Basanta, E Bello-López, et al. 2022. {``Genomic
Characterization of Antifungal {Acinetobacter} Bacteria Isolated from
the Skin of the Frogs {Agalychnis} Callidryas and {Craugastor}
Fitzingeri.''} \emph{FEMS Microbiology Ecology} 98 (12): fiac126.
\url{https://doi.org/10.1093/femsec/fiac126}.

\bibitem[\citeproctext]{ref-chambers_statistical_1992}
Chambers, J. M., and T. J. Hastie. 1992. \emph{Statistical {Models} in
{S}}. Wadsworth \& Brooks/Cole.

\bibitem[\citeproctext]{ref-chao_unifying_2014}
Chao, Anne, Chun-Huo Chiu, and Lou Jost. 2014. {``Unifying Species
Diversity, Phylogenetic Diversity, Functional Diversity, and Related
Similarity and Differentiation Measures Through {Hill} Numbers.''}
\emph{Annual Review of Ecology, Evolution, and Systematics} 45 (1):
297--324. \url{https://doi.org/10.1146/annurev-ecolsys-120213-091540}.

\bibitem[\citeproctext]{ref-chaumeil_gtdb-tk_2022}
Chaumeil, Pierre-Alain, Aaron J Mussig, Philip Hugenholtz, and Donovan H
Parks. 2022. {``{GTDB}-{Tk} V2: Memory Friendly Classification with the
Genome Taxonomy Database.''} \emph{Bioinformatics} 38 (23): 5315--16.
\url{https://doi.org/10.1093/bioinformatics/btac672}.

\bibitem[\citeproctext]{ref-chen_fastp_2018}
Chen, Shifu, Yanqing Zhou, Yaru Chen, and Jia Gu. 2018. {``Fastp: An
Ultra-Fast All-in-One {FASTQ} Preprocessor.''} \emph{Bioinformatics} 34
(17): i884--90. \url{https://doi.org/10.1093/bioinformatics/bty560}.

\bibitem[\citeproctext]{ref-chklovski_checkm2_2023}
Chklovski, Alex, Donovan H Parks, Ben J Woodcroft, and Gene W Tyson.
2023. {``{CheckM2}: A Rapid, Scalable and Accurate Tool for Assessing
Microbial Genome Quality Using Machine Learning.''} \emph{Nature
Methods} 20 (8): 1203--12.
\url{https://doi.org/10.1038/s41592-023-01940-w}.

\bibitem[\citeproctext]{ref-clavel_worldwide_2011}
Clavel, Joanne, Romain Julliard, and Vincent Devictor. 2011.
{``Worldwide Decline of Specialist Species: Toward a Global Functional
Homogenization?''} \emph{Frontiers in Ecology and the Environment} 9
(4): 222--28. \url{https://doi.org/10.1890/080216}.

\bibitem[\citeproctext]{ref-coban_soil_2022}
Coban, Oksana, Gerlinde B De Deyn, and Martine van der Ploeg. 2022.
{``Soil Microbiota as Game-Changers in Restoration of Degraded Lands.''}
\emph{Science} 375 (6584): abe0725.
\url{https://doi.org/10.1126/science.abe0725}.

\bibitem[\citeproctext]{ref-cordeiro_amphibian_2024}
{Cordeiro, Isabella Ferreira, Camila Gracyelle de Carvalho Lemes,
Angélica Bianchini Sanchez, et al.} 2024. {``Amphibian Tolerance to
Arsenic: Microbiome-Mediated Insights.''} \emph{Scientific Reports} 14
(1): 10193. \url{https://doi.org/10.1038/s41598-024-60879-w}.

\bibitem[\citeproctext]{ref-cox_global_2022}
Cox, Neil, Bruce E. Young, Philip Bowles, et al. 2022. {``A Global
Reptile Assessment Highlights Shared Conservation Needs of Tetrapods.''}
\emph{Nature} 605 (7909): 285--90.
\url{https://doi.org/10.1038/s41586-022-04664-7}.

\bibitem[\citeproctext]{ref-crawshaw_widespread_2022}
Crawshaw, Lauren, Tore Buchanan, Leonard Shirose, et al. 2022.
{``Widespread Occurrence of {Batrachochytrium} Dendrobatidis in
{Ontario}, {Canada}, and Predicted Habitat Suitability for the Emerging
{Batrachochytrium} Salamandrivorans.''} \emph{Ecology and Evolution} 12
(4): e8798. \url{https://doi.org/10.1002/ece3.8798}.

\bibitem[\citeproctext]{ref-danecek_twelve_2021}
{Danecek, Petr, James K Bonfield, Jennifer Liddle, et al.} 2021.
{``Twelve Years of {SAMtools} and {BCFtools}.''} \emph{Gigascience} 10
(2): giab008. \url{https://doi.org/10.1093/gigascience/giab008}.

\bibitem[\citeproctext]{ref-demori_peptides_2019}
{Demori, Ilaria, Zeinab El Rashed, Viola Corradino, et al.} 2019.
{``Peptides for Skin Protection and Healing in Amphibians.''}
\emph{Molecules} 24 (2): 347.
\url{https://doi.org/10.3390/molecules24020347}.

\bibitem[\citeproctext]{ref-dinis_dry-climate_2025}
Dinis, Marco, Fernando Martínez-Freiría, Wouter Beukema, Sergé Bogaerts,
and Guillermo Velo-Antón. 2025. {``The {Dry}-{Climate} {Hypothesis}:
{Identifying} the {Environmental} {Drivers} of {Terrestrial}
{Viviparous} {Salamanders}.''} \emph{Journal of Biogeography} 52 (11):
e70046. \url{https://doi.org/10.1111/jbi.70046}.

\bibitem[\citeproctext]{ref-dong_rice_2024}
Dong, Wen-jing, Ming-da Xu, Xue-wen Yang, et al. 2024. {``Rice Straw Ash
and Amphibian Health: {A} Deep Dive into Microbiota Changes and
Potential Ecological Consequences.''} \emph{Science of The Total
Environment} 926: 171651.
\url{https://doi.org/10.1016/j.scitotenv.2024.171651}.

\bibitem[\citeproctext]{ref-durazzi_comparison_2021}
Durazzi, Francesco, Claudia Sala, Gastone Castellani, Gerardo Manfreda,
Daniel Remondini, and Alessandra De Cesare. 2021. {``Comparison Between
{16S} {rRNA} and Shotgun Sequencing Data for the Taxonomic
Characterization of the Gut Microbiota.''} \emph{Scientific Reports} 11
(1): 3030. \url{https://doi.org/10.1038/s41598-021-82726-y}.

\bibitem[\citeproctext]{ref-eisenhofer_quantifying_2024}
Eisenhofer, Raphael, Antton Alberdi, and Ben J Woodcroft. 2024.
{``Quantifying Microbial {DNA} in Metagenomes Improves Microbial Trait
Estimation.''} \emph{ISME Communications} 4 (1): ycae111.
\url{https://doi.org/10.1093/ismeco/ycae111}.

\bibitem[\citeproctext]{ref-eisenhofer_impact_2023}
Eisenhofer, Raphael, Iñaki Odriozola, and Antton Alberdi. 2023.
{``Impact of Microbial Genome Completeness on Metagenomic Functional
Inference.''} \emph{ISME Communications} 3 (1): 12.
\url{https://doi.org/10.1038/s43705-023-00221-z}.

\bibitem[\citeproctext]{ref-faulhaber_detection_2025}
Faulhaber, Marline M., Florence Tardy, Franziska Saul, Elisabeth Müller,
Michael Pees, and Rachel E. Marschang. 2025. {``Detection of
{Mycoplasma} Spp. From Snakes from Five Different Families.''} \emph{BMC
Veterinary Research} 21 (1): 38.
\url{https://doi.org/10.1186/s12917-025-04487-4}.

\bibitem[\citeproctext]{ref-ferreira_lipopolysaccharide_2021}
Ferreira, L F, P G Garcia Neto, S C M Titon, et al. 2021.
{``Lipopolysaccharide {Regulates} {Pro}- and {Anti}-{Inflammatory}
{Cytokines}, {Corticosterone}, and {Melatonin} in {Toads}.''}
\emph{Integrative Organismal Biology} 3 (1): obab025.
\url{https://doi.org/10.1093/iob/obab025}.

\bibitem[\citeproctext]{ref-finn_pfam_2016}
{Finn, Robert D, Penelope Coggill, Ruth Y Eberhardt, et al.} 2016.
{``The {Pfam} Protein Families Database: Towards a More Sustainable
Future.''} \emph{Nucleic Acids Research} 44 (D1): D279--85.
\url{https://doi.org/10.1093/nar/gkv1344}.

\bibitem[\citeproctext]{ref-fisher_chytrid_2020}
Fisher, Matthew C, and Trenton WJ Garner. 2020. {``Chytrid Fungi and
Global Amphibian Declines.''} \emph{Nature Reviews Microbiology} 18 (6):
332--43. \url{https://doi.org/10.1038/s41579-020-0335-x}.

\bibitem[\citeproctext]{ref-friedman_analysis_2018}
Friedman, Mendel. 2018. {``Analysis, {Nutrition}, and {Health}
{Benefits} of {Tryptophan}.''} \emph{International Journal of Tryptophan
Research} 11 (January): 1178646918802282.
\url{https://doi.org/10.1177/1178646918802282}.

\bibitem[\citeproctext]{ref-gloor_microbiome_2017}
Gloor, Gregory B, Jean M Macklaim, Vera Pawlowsky-Glahn, and Juan J
Egozcue. 2017. {``Microbiome Datasets Are Compositional: And This Is Not
Optional.''} \emph{Frontiers in Microbiology} 8: 2224.
\url{https://doi.org/10.3389/fmicb.2017.02224}.

\bibitem[\citeproctext]{ref-gloor_compositional_2016}
Gloor, Gregory B, and Gregor Reid. 2016. {``Compositional Analysis: A
Valid Approach to Analyze Microbiome High-Throughput Sequencing Data.''}
\emph{Canadian Journal of Microbiology} 62 (8): 692--703.
\url{https://doi.org/10.1139/cjm-2015-0821}.

\bibitem[\citeproctext]{ref-gonzalez-serrano_amphibian_2025}
{González-Serrano, Francisco, Yordan J Romero-Contreras, Alberto H Orta,
et al.} 2025. {``Amphibian Skin Bacteria Contain a Wide Repertoire of
Genes Linked to Their Antifungal Capacities.''} \emph{World Journal of
Microbiology and Biotechnology} 41 (3): 78.
\url{https://doi.org/10.1007/s11274-025-04292-z}.

\bibitem[\citeproctext]{ref-guilherme_looking_2022}
Guilherme, Filipa, Eva García Moreno, José Alberto Gonçalves, Miguel A
Carretero, and Paulo Farinha-Marques. 2022. {``Looking Closer at the
Patterns of Land Cover in the {City} of {Porto}, {Portugal}, Between
1947 and 2019---{A} Contribution for the Integration of Ecological Data
in Spatial Planning.''} \emph{Land} 11 (10): 1828.
\url{https://doi.org/10.3390/land11101828}.

\bibitem[\citeproctext]{ref-gurevich_quast_2013}
Gurevich, Alexey, Vladislav Saveliev, Nikolay Vyahhi, and Glenn Tesler.
2013. {``{QUAST}: Quality Assessment Tool for Genome Assemblies.''}
\emph{Bioinformatics} 29 (8): 1072--75.
\url{https://doi.org/10.1093/bioinformatics/btt086}.

\bibitem[\citeproctext]{ref-hamer_amphibian_2008}
Hamer, Andrew J, and Mark J McDonnell. 2008. {``Amphibian Ecology and
Conservation in the Urbanising World: A Review.''} \emph{Biological
Conservation} 141 (10): 2432--49.
\url{https://doi.org/10.1016/j.biocon.2008.07.020}.

\bibitem[\citeproctext]{ref-hernandez_contrasting_2024}
Hernández, Mauricio, Jorge Langa, Ostaizka Aizpurua, Yendi E
Navarro-Noya, and Antton Alberdi. 2024. \emph{Contrasting Recovery of
Metagenome-Assembled Genomes and Derived Microbial Communities from
Lizard Fecal and Cloacal Samples}.
\url{https://doi.org/10.21203/rs.3.rs-5287039/v1}.

\bibitem[\citeproctext]{ref-hernandez-gomez_invasive_2020}
Hernández-Gómez, Obed, Allison Q Byrne, Alex R Gunderson, et al. 2020.
{``Invasive Vegetation Affects Amphibian Skin Microbiota and Body
Condition.''} \emph{PeerJ} 8: e8549.
\url{https://doi.org/10.7717/peerj.8549}.

\bibitem[\citeproctext]{ref-hernandez-gomez_organismal_2023}
Hernández-Gómez, Obed, and Jessica Hua. 2023. {``From the Organismal to
Biosphere Levels: Environmental Impacts on the Amphibian Microbiota.''}
\emph{FEMS Microbiology Reviews} 47 (1): fuad002.
\url{https://doi.org/doi.org/10.1093/femsre/fuad002}.

\bibitem[\citeproctext]{ref-hill_diversity_1973}
Hill, Mark O. 1973. {``Diversity and Evenness: A Unifying Notation and
Its Consequences.''} \emph{Ecology} 54 (2): 427--32.
\url{https://doi.org/10.2307/1934352}.

\bibitem[\citeproctext]{ref-hoyt_bacteria_2015}
Hoyt, Joseph R, Tina L Cheng, Kate E Langwig, Mallory M Hee, Winifred F
Frick, and A Marm Kilpatrick. 2015. {``Bacteria Isolated from Bats
Inhibit the Growth of {Pseudogymnoascus} Destructans, the Causative
Agent of White-Nose Syndrome.''} \emph{PLoS One} 10 (4): e0121329.
\url{https://doi.org/10.1371/journal.pone.0121329}.

\bibitem[\citeproctext]{ref-hu_impervious_2018}
Hu, Yinhong, Xiaolin Dou, Juanyong Li, and Feng Li. 2018. {``Impervious
Surfaces Alter Soil Bacterial Communities in Urban Areas: A Case Study
in {Beijing}, {China}.''} \emph{Frontiers in Microbiology} 9: 226.
\url{https://doi.org/10.3389/fmicb.2018.00226}.

\bibitem[\citeproctext]{ref-iacobazzi_citrate_2014}
Iacobazzi, Vito, and Vittoria Infantino. 2014. \emph{Citrate -- New
Functions for an Old Metabolite}. Biological {Chemistry}, vol. 395 (4):
387--99. \url{https://doi.org/10.1515/hsz-2013-0271}.

\bibitem[\citeproctext]{ref-iglesias-carrasco_urban_2017}
Iglesias-Carrasco, Maider, José Martín, and Carlos Cabido. 2017.
{``Urban Habitats Can Affect Body Size and Body Condition but Not Immune
Response in Amphibians.''} \emph{Urban Ecosystems} 20: 1331--38.
\url{https://doi.org/10.1007/s11252-017-0685-y}.

\bibitem[\citeproctext]{ref-jani_amphibian_2021}
Jani, Andrea J, Jessie Bushell, Cédric G Arisdakessian, et al. 2021.
{``The Amphibian Microbiome Exhibits Poor Resilience Following
Pathogen-Induced Disturbance.''} \emph{The ISME Journal} 15 (6):
1628--40. \url{https://doi.org/10.1038/s41396-020-00875-w}.

\bibitem[\citeproctext]{ref-jennette_variation_2019}
Jennette, Matthew A, Joel W Snodgrass, and Don C Forester. 2019.
{``Variation in Age, Body Size, and Reproductive Traits Among Urban and
Rural Amphibian Populations.''} \emph{Urban Ecosystems} 22 (1): 137--47.
\url{https://doi.org/10.1007/s11252-018-0801-7}.

\bibitem[\citeproctext]{ref-jimenez_fungicide_2021}
Jiménez, Randall R, Gilbert Alvarado, Clemens Ruepert, Erick Ballestero,
and Simone Sommer. 2021. {``The Fungicide Chlorothalonil Changes the
Amphibian Skin Microbiome: A Potential Factor Disrupting a Host
Disease-Protective Trait.''} \emph{Applied Microbiology} 1 (1): 26--37.
\url{https://doi.org/10.3390/applmicrobiol1010004}.

\bibitem[\citeproctext]{ref-kanehisa_kegg_2017}
Kanehisa, Minoru, Miho Furumichi, Mao Tanabe, Yoko Sato, and Kanae
Morishima. 2017. {``{KEGG}: New Perspectives on Genomes, Pathways,
Diseases and Drugs.''} \emph{Nucleic Acids Research} 45 (D1): D353--61.
\url{https://doi.org/10.1093/nar/gkw1092}.

\bibitem[\citeproctext]{ref-kang_metabat_2019}
Kang, Dongwan D, Feng Li, Edward Kirton, et al. 2019. {``{MetaBAT} 2: An
Adaptive Binning Algorithm for Robust and Efficient Genome
Reconstruction from Metagenome Assemblies.''} \emph{PeerJ} 7: e7359.
\url{https://doi.org/10.7717/peerj.7359}.

\bibitem[\citeproctext]{ref-kaufmann_adaption_2016}
Kaufmann, Katharina, and Peter Dohmen. 2016. {``Adaption of a Dermal in
Vitro Method to Investigate the Uptake of Chemicals Across Amphibian
Skin.''} \emph{Environmental Sciences Europe} 28 (1): 10.
\url{https://doi.org/10.1186/s12302-016-0080-y}.

\bibitem[\citeproctext]{ref-kosch_comparative_2025}
Kosch, Tiffany A., Andrew J. Crawford, Rachel Lockridge Mueller, et al.
2025. {``Comparative Analysis of Amphibian Genomes: {An} Emerging
Resource for Basic and Applied Research.''} \emph{Molecular Ecology
Resources} 25 (1): e14025.
\url{https://doi.org/10.1111/1755-0998.14025}.

\bibitem[\citeproctext]{ref-kosch_amphibian_2024}
{Kosch, Tiffany A, María Torres-Sánchez, H Christoph Liedtke, et al.}
2024. {``The {Amphibian} {Genomics} {Consortium}: Advancing Genomic and
Genetic Resources for Amphibian Research and Conservation.''} \emph{BMC
Genomics} 25 (1): 1025.
\url{https://doi.org/10.1186/s12864-024-10899-7}.

\bibitem[\citeproctext]{ref-kueneman_effects_2022}
{Kueneman, Jordan G, Molly C Bletz, Matthew Becker, et al.} 2022.
{``Effects of Captivity and Rewilding on Amphibian Skin Microbiomes.''}
\emph{Biological Conservation} 271: 109576.
\url{https://doi.org/10.1016/j.biocon.2022.109576}.

\bibitem[\citeproctext]{ref-kueneman_amphibian_2014}
Kueneman, Jordan G, Laura Wegener Parfrey, Douglas C Woodhams, Holly M
Archer, Rob Knight, and Valerie J McKenzie. 2014. {``The Amphibian
Skin-Associated Microbiome Across Species, Space and Life History
Stages.''} \emph{Molecular Ecology} 23 (6): 1238--50.
\url{https://doi.org/10.1111/mec.12510}.

\bibitem[\citeproctext]{ref-kurtz_sparse_2015}
Kurtz, Zachary D, Christian L Müller, Emily R Miraldi, Dan R Littman,
Martin J Blaser, and Richard A Bonneau. 2015. {``Sparse and
Compositionally Robust Inference of Microbial Ecological Networks.''}
\emph{PLoS Computational Biology} 11 (5): e1004226.
\url{https://doi.org/10.1371/journal.pcbi.1004226}.

\bibitem[\citeproctext]{ref-langmead_fast_2012}
Langmead, Ben, and Steven L Salzberg. 2012. {``Fast Gapped-Read
Alignment with {Bowtie} 2.''} \emph{Nature Methods} 9 (4): 357--59.
\url{https://doi.org/10.1038/nmeth.1923}.

\bibitem[\citeproctext]{ref-lauritsen_robust_2025}
{Lauritsen, Jonas G, Christian Carøe, Nanna Gaun, et al.} 2025.
{``Robust, {Open}-{Source} and {Automation}-{Friendly} {DNA}
{Extraction} {Protocol} for {Hologenomic} {Research}.''} \emph{Molecular
Ecology Resources} 25 (8): e70042.
\url{https://doi.org/10.1111/1755-0998.70042}.

\bibitem[\citeproctext]{ref-le_sage_preparatory_2021}
Le Sage, Emily H, Brandon C LaBumbard, Laura K Reinert, et al. 2021.
{``Preparatory Immunity: {Seasonality} of Mucosal Skin Defences and
{Batrachochytrium} Infections in {Southern} Leopard Frogs.''}
\emph{Journal of Animal Ecology} 90 (2): 542--54.
\url{https://doi.org/10.1111/1365-2656.13386}.

\bibitem[\citeproctext]{ref-leonard_global_2024}
{Leonard, Aoife, Javier Abalos, Titus Adhola, et al.} 2024. {``A Global
Initiative for Ecological and Evolutionary Hologenomics.''} \emph{Trends
in Ecology \& Evolution}, ahead of print.
\url{https://doi.org/10.1016/j.tree.2024.03.005}.

\bibitem[\citeproctext]{ref-li_megahit_2016}
Li, Dinghua, Ruibang Luo, Chi-Man Liu, et al. 2016. {``{MEGAHIT} V1. 0:
A Fast and Scalable Metagenome Assembler Driven by Advanced
Methodologies and Community Practices.''} \emph{Methods} 102: 3--11.
\url{https://doi.org/10.1016/j.ymeth.2016.02.020}.

\bibitem[\citeproctext]{ref-li_damage_2022}
Li, Minghao, Jincheng Wang, Huiling Deng, et al. 2022. {``The {Damage}
of the {Crayfish} ({Procambarus} {Clarkii}) {Digestive} {Organs}
{Caused} by {Citrobacter} {Freundii} {Is} {Associated} {With} the
{Disturbance} of {Intestinal} {Microbiota} and {Disruption} of
{Intestinal}-{Liver} {Axis} {Homeostasis}.''} \emph{Frontiers in
Cellular and Infection Microbiology} Volume 12 - 2022.
\url{https://www.frontiersin.org/journals/cellular-and-infection-microbiology/articles/10.3389/fcimb.2022.940576}.

\bibitem[\citeproctext]{ref-libro_novo_2023}
Libro, Pietro, Andrea Chiocchio, Erika De Rysky, et al. 2023. {``De Novo
Transcriptome Assembly and Annotation for Gene Discovery in {Salamandra}
Salamandra at the Larval Stage.''} \emph{Scientific Data} 10 (1): 330.
\url{https://doi.org/10.1038/s41597-023-02217-9}.

\bibitem[\citeproctext]{ref-liedtke_macroevolutionary_2018}
Liedtke, H Christoph, David J Gower, Mark Wilkinson, and Ivan
Gomez-Mestre. 2018. {``Macroevolutionary Shift in the Size of Amphibian
Genomes and the Role of Life History and Climate.''} \emph{Nature
Ecology \& Evolution} 2 (11): 1792--99.
\url{https://doi.org/10.1038/s41559-018-0674-4}.

\bibitem[\citeproctext]{ref-lin_analysis_2020}
Lin, Huang, and Shyamal Das Peddada. 2020. {``Analysis of Compositions
of Microbiomes with Bias Correction.''} \emph{Nature Communications} 11
(1): 3514. \url{https://doi.org/10.1038/s41467-020-17041-7}.

\bibitem[\citeproctext]{ref-liu_microbial_2025}
Liu, Pengfei, Shizhi Wen, Shanshan Zhu, Xi Hu, and Yamin Wang. 2025.
{``Microbial {Degradation} of {Soil} {Organic} {Pollutants}:
{Mechanisms}, {Challenges}, and {Advances} in {Forest} {Ecosystem}
{Management}.''} \emph{Processes} 13 (3): 916.
\url{https://doi.org/10.3390/pr13030916}.

\bibitem[\citeproctext]{ref-longcore_batrachochytrium_1999}
Longcore, Joyce E, Allan P Pessier, and Donald K Nichols. 1999.
{``Batrachochytrium Dendrobatidis Gen. Et Sp. Nov., a Chytrid Pathogenic
to Amphibians.''} \emph{Mycologia} 91 (2): 219--27.
\url{https://doi.org/10.1080/00275514.1999.12061011}.

\bibitem[\citeproctext]{ref-loudon_microbial_2013}
Loudon, Andrew H, Douglas C Woodhams, Laura Wegener Parfrey, et al.
2013. {``Microbial Community Dynamics and Effect of Environmental
Microbial Reservoirs on Red-Backed Salamanders ({Plethodon}
Cinereus).''} \emph{The ISME Journal} 8 (4): 830--40.
\url{https://doi.org/10.1038/ismej.2013.200}.

\bibitem[\citeproctext]{ref-lourenco_trapped_2017}
Lourenço, André, David Álvarez, Ian J Wang, and Guillermo Velo-Antón.
2017. {``Trapped Within the City: {Integrating} Demography, Time Since
Isolation and Population-Specific Traits to Assess the Genetic Effects
of Urbanization.''} \emph{Molecular Ecology} 26 (6): 1498--514.
\url{https://doi.org/10.1111/mec.14019}.

\bibitem[\citeproctext]{ref-luddecke_salamanders_2018}
Lüddecke, Tim, Stefan Schulz, Sebastian Steinfartz, and Miguel Vences.
2018. {``A Salamander's Toxic Arsenal: Review of Skin Poison Diversity
and Function in True Salamanders, Genus {Salamandra}.''} \emph{The
Science of Nature} 105 (9): 56.
\url{https://doi.org/10.1007/s00114-018-1579-4}.

\bibitem[\citeproctext]{ref-luedtke_ongoing_2023}
{Luedtke, Jennifer A, Janice Chanson, Kelsey Neam, et al.} 2023.
{``Ongoing Declines for the World's Amphibians in the Face of Emerging
Threats.''} \emph{Nature} 622 (7982): 308--14.
\url{https://doi.org/10.1038/s41586-023-06578-4}.

\bibitem[\citeproctext]{ref-macdonald_meta-analysis_2023}
Macdonald, Kristina J, Don A Driscoll, Kimberley J Macdonald, Bronwyn
Hradsky, and Tim S Doherty. 2023. {``Meta-Analysis Reveals Impacts of
Disturbance on Reptile and Amphibian Body Condition.''} \emph{Global
Change Biology} 29 (17): 4949--65.
\url{https://doi.org/10.1111/gcb.16852}.

\bibitem[\citeproctext]{ref-madureira_green_2011}
Madureira, Helena, Teresa Andresen, and Ana Monteiro. 2011. {``Green
Structure and Planning Evolution in {Porto}.''} \emph{Urban Forestry \&
Urban Greening} 10 (2): 141--49.
\url{https://doi.org/10.1016/j.ufug.2010.12.004}.

\bibitem[\citeproctext]{ref-mandal_analysis_2015}
Mandal, Siddhartha, Will Van Treuren, Richard A White, Merete Eggesbø,
Rob Knight, and Shyamal D Peddada. 2015. {``Analysis of Composition of
Microbiomes: A Novel Method for Studying Microbial Composition.''}
\emph{Microbial Ecology in Health and Disease} 26 (1): 27663.
\url{https://doi.org/10.3402/mehd.v26.27663}.

\bibitem[\citeproctext]{ref-marques_coi_2022}
Marques, Adam J. D., Vanessa A. Mata, and Guillermo Velo-Antón. 2022.
{``{COI} {Metabarcoding} {Provides} {Insights} into the {Highly}
{Diverse} {Diet} of a {Generalist} {Salamander}, {Salamandra} Salamandra
({Caudata}: {Salamandridae}).''} \emph{Diversity} 14 (2): 89.
\url{https://doi.org/10.3390/d14020089}.

\bibitem[\citeproctext]{ref-martel_batrachochytrium_2013}
{Martel, An, Annemarieke Spitzen-van der Sluijs, Mark Blooi, et al.}
2013. {``Batrachochytrium Salamandrivorans Sp. Nov. Causes Lethal
Chytridiomycosis in Amphibians.''} \emph{Proceedings of the National
Academy of Sciences} 110 (38): 15325--29.
\url{https://doi.org/10.1073/pnas.1307356110}.

\bibitem[\citeproctext]{ref-mayeux_pathobiology_1997}
Mayeux, Philip R. 1997. {``Pathobiology of {Lipolysaccharide}.''}
\emph{Journal of Toxicology and Environmental Health} 51 (5): 415--35.
\url{https://doi.org/10.1080/00984109708984034}.

\bibitem[\citeproctext]{ref-mcgarvey_uniprot_2019}
McGarvey, Peter B, Andrew Nightingale, Jie Luo, et al. 2019.
{``{UniProt} Genomic Mapping for Deciphering Functional Effects of
Missense Variants.''} \emph{Human Mutation} 40 (6): 694--705.
\url{https://doi.org/10.1002/humu.23738}.

\bibitem[\citeproctext]{ref-mckinney_urbanization_2006}
McKinney, Michael L. 2006. {``Urbanization as a Major Cause of Biotic
Homogenization.''} \emph{Biological Conservation} 127 (3): 247--60.
\url{https://doi.org/10.1016/j.biocon.2005.09.005}.

\bibitem[\citeproctext]{ref-mclellan_microbiome_2015}
McLellan, Sandra L, Jenny C Fisher, and Ryan J Newton. 2015. {``The
Microbiome of Urban Waters.''} \emph{International Microbiology: The
Official Journal of the Spanish Society for Microbiology} 18 (3):
141--49. \url{https://doi.org/10.2436/20.1501.01.244}.

\bibitem[\citeproctext]{ref-mcnabb_effects_2019}
McNabb, Micaiah C, Christa Marie Saelinger, Melody Danley, and Robin
Lewis Cooper. 2019. {``Effects of {Bacterial} {Endotoxin}
{Lipopolysaccharides} ({LPS}) on {Synaptic} {Transmission} at
{Neuromuscular} {Junction} in an {Amphibian}.''} \emph{Int. J. Zool.
Res} 15: 38--42. \url{https://doi.org/10.3923/ijzr.2019.38.42}.

\bibitem[\citeproctext]{ref-morell_frogs_1999}
Morell, Virginia. 1999. {``Frogs: {Canaries} in the {Hot} {Zone}?''}
\emph{Science} 284 (5415): 729--29.
\url{https://doi.org/10.1126/science.284.5415.729}.

\bibitem[\citeproctext]{ref-muletz_soil_2012}
Muletz, Carly R, Jillian M Myers, Rickie J Domangue, James B Herrick,
and Reid N Harris. 2012. {``Soil Bioaugmentation with Amphibian
Cutaneous Bacteria Protects Amphibian Hosts from Infection by
{Batrachochytrium} Dendrobatidis.''} \emph{Biological Conservation} 152:
119--26. \url{https://doi.org/10.1016/j.biocon.2012.03.022}.

\bibitem[\citeproctext]{ref-muletz_wolz_effects_2018}
Muletz Wolz, Carly R, Stephanie A Yarwood, Evan H Campbell Grant, Robert
C Fleischer, and Karen R Lips. 2018. {``Effects of Host Species and
Environment on the Skin Microbiome of {Plethodontid} Salamanders.''}
\emph{Journal of Animal Ecology} 87 (2): 341--53.
\url{https://doi.org/10.1111/1365-2656.12726}.

\bibitem[\citeproctext]{ref-muletz-wolz_inhibition_2017}
Muletz-Wolz, Carly R, Jose G Almario, Samuel E Barnett, et al. 2017.
{``Inhibition of Fungal Pathogens Across Genotypes and Temperatures by
Amphibian Skin Bacteria.''} \emph{Frontiers in Microbiology} 8: 1551.
\url{https://doi.org/10.3389/fmicb.2017.01551}.

\bibitem[\citeproctext]{ref-murray_benchtop_2015}
Murray, Dáithí C, Megan L Coghlan, and Michael Bunce. 2015. {``From
Benchtop to Desktop: Important Considerations When Designing Amplicon
Sequencing Workflows.''} \emph{PloS One} 10 (4): e0124671.
\url{https://doi.org/10.1371/journal.pone.0124671}.

\bibitem[\citeproctext]{ref-nearing_identifying_2021}
Nearing, Jacob T, André M Comeau, and Morgan GI Langille. 2021.
{``Identifying Biases and Their Potential Solutions in Human Microbiome
Studies.''} \emph{Microbiome} 9 (1): 113.
\url{https://doi.org/10.1186/s40168-021-01059-0}.

\bibitem[\citeproctext]{ref-norris_ghosts_2007}
Norris, Scott. 2007. {``Ghosts in Our Midst: Coming to Terms with
Amphibian Extinctions.''} \emph{BioScience} 57 (4): 311--16.
\url{https://doi.org/10.1641/B570403}.

\bibitem[\citeproctext]{ref-ohanlon_recent_2018}
{O'Hanlon, Simon J, Adrien Rieux, Rhys A Farrer, et al.} 2018. {``Recent
{Asian} Origin of Chytrid Fungi Causing Global Amphibian Declines.''}
\emph{Science} 360 (6389): 621--27.
\url{https://doi.org/10.1126/science.aar1965}.

\bibitem[\citeproctext]{ref-ohmer_skin_2015}
Ohmer, Michel E. B., Rebecca L. Cramp, Craig R. White, and Craig E.
Franklin. 2015. {``Skin Sloughing Rate Increases with Chytrid Fungus
Infection Load in a Susceptible Amphibian.''} \emph{Functional Ecology}
29 (5): 674--82. \url{https://doi.org/10.1111/1365-2435.12370}.

\bibitem[\citeproctext]{ref-oksanen_package_2013}
{Oksanen, Jari, F Guillaume Blanchet, Roeland Kindt, et al.} 2013.
{``Package {`Vegan'}.''} \emph{Community Ecology Package, Version} 2
(9): 1--295.

\bibitem[\citeproctext]{ref-olm_drep_2017}
Olm, Matthew R, Christopher T Brown, Brandon Brooks, and Jillian F
Banfield. 2017. {``{dRep}: A Tool for Fast and Accurate Genomic
Comparisons That Enables Improved Genome Recovery from Metagenomes
Through de-Replication.''} \emph{The ISME Journal} 11 (12): 2864--68.
\url{https://doi.org/10.1038/ismej.2017.126}.

\bibitem[\citeproctext]{ref-paranhos_sepsis_2024}
Paranhos, Gabriel F., Hodias S. Oliveira Filho, José L. C. Duarte, et
al. 2024. {``Sepsis by {Plesiomonas} Shigelloides, {Citrobacter}
Freundii and {Aeromonas} Jandaei in a Green Iguana ({Iguana} Iguana).''}
\emph{Pesquisa Veterinária Brasileira} 44.

\bibitem[\citeproctext]{ref-parks_gtdb_2022}
Parks, Donovan H, Maria Chuvochina, Christian Rinke, Aaron J Mussig,
Pierre-Alain Chaumeil, and Philip Hugenholtz. 2022. {``{GTDB}: An
Ongoing Census of Bacterial and Archaeal Diversity Through a
Phylogenetically Consistent, Rank Normalized and Complete Genome-Based
Taxonomy.''} \emph{Nucleic Acids Research} 50 (D1): D785--94.
\url{https://doi.org/10.1093/nar/gkab776}.

\bibitem[\citeproctext]{ref-pereira_skin_2018}
Pereira, Kenzie E, Brian I Crother, David M Sever, et al. 2018. {``Skin
Glands of an Aquatic Salamander Vary in Size and Distribution and
Release Antimicrobial Secretions Effective Against Chytrid Fungal
Pathogens.''} \emph{Journal of Experimental Biology} 221 (14):
jeb183707. \url{https://doi.org/10.1242/jeb.183707}.

\bibitem[\citeproctext]{ref-pereira-marques_impact_2019}
Pereira-Marques, Joana, Anne Hout, Rui M Ferreira, et al. 2019.
{``Impact of Host {DNA} and Sequencing Depth on the Taxonomic Resolution
of Whole Metagenome Sequencing for Microbiome Analysis.''}
\emph{Frontiers in Microbiology} 10: 1277.
\url{https://doi.org/10.3389/fmicb.2019.01277}.

\bibitem[\citeproctext]{ref-pietroni_hologenomic_2025}
Pietroni, Carlotta, Nanna Gaun, Aoife Leonard, et al. 2025.
{``Hologenomic Data Generation and Analysis in Wild Vertebrates.''}
\emph{Methods in Ecology and Evolution} 16 (1): 97--107.
\url{https://doi.org/10.1111/2041-210X.14456}.

\bibitem[\citeproctext]{ref-pruesse_silva_2007}
Pruesse, Elmar, Christian Quast, Katrin Knittel, et al. 2007.
{``{SILVA}: A Comprehensive Online Resource for Quality Checked and
Aligned Ribosomal {RNA} Sequence Data Compatible with {ARB}.''}
\emph{Nucleic Acids Research} 35 (21): 7188--96.
\url{https://doi.org/10.1093/nar/gkm864}.

\bibitem[\citeproctext]{ref-purucker_estimating_2023}
Purucker, S Thomas, Marcía N Snyder, Donna A Glinski, et al. 2023.
{``Estimating Dermal Contact Soil Exposure for Amphibians.''}
\emph{Integrated Environmental Assessment and Management} 19 (1): 9--16.
\url{https://doi.org/10.1002/ieam.4619}.

\bibitem[\citeproctext]{ref-quaranta_why_2009}
Quaranta, Angelo, Vito Bellantuono, Giuseppe Cassano, and Claudio Lippe.
2009. {``Why Amphibians Are More Sensitive Than Mammals to
Xenobiotics.''} \emph{PLoS One} 4 (11): e7699.
\url{https://doi.org/10.1371/journal.pone.0007699}.

\bibitem[\citeproctext]{ref-quast_silva_2013}
Quast, Christian, Elmar Pruesse, Pelin Yilmaz, et al. 2013. {``The
{SILVA} Ribosomal {RNA} Gene Database Project: Improved Data Processing
and Web-Based Tools.''} \emph{Nucleic Acids Research} 41 (D1): D590--96.
\url{https://doi.org/10.1093/nar/gks1219}.

\bibitem[\citeproctext]{ref-quince_shotgun_2017}
Quince, Christopher, Alan W Walker, Jared T Simpson, Nicholas J Loman,
and Nicola Segata. 2017. {``Shotgun Metagenomics, from Sampling to
Analysis.''} \emph{Nature Biotechnology} 35 (9): 833--44.
\url{https://doi.org/10.1038/nbt.3935}.

\bibitem[\citeproctext]{ref-raetz_lipopolysaccharide_2002}
Raetz, Christian R. H., and Chris Whitfield. 2002. {``Lipopolysaccharide
{Endotoxins}.''} \emph{Annual Review of Biochemistry} 71 (Volume 71,
2002): 635--700.
https://doi.org/\url{https://doi.org/10.1146/annurev.biochem.71.110601.135414}.

\bibitem[\citeproctext]{ref-rawlings_merops_1999}
Rawlings, Neil D, and Alan J Barrett. 1999. {``{MEROPS}: The Peptidase
Database.''} \emph{Nucleic Acids Research} 27 (1): 325--31.
\url{https://doi.org/10.1093/nar/gkp971}.

\bibitem[\citeproctext]{ref-rollins-smith_batrachochytrium_2021}
Rollins-Smith, Louise A, and Emily H Le Sage. 2021. {``Batrachochytrium
Fungi: Stealth Invaders in Amphibian Skin.''} \emph{Current Opinion in
Microbiology} 61: 124--32.
\url{https://doi.org/10.1016/j.mib.2021.04.002}.

\bibitem[\citeproctext]{ref-ross_skin_2019}
Ross, Ashley A., Aline Rodrigues Hoffmann, and Josh D. Neufeld. 2019.
{``The Skin Microbiome of Vertebrates.''} \emph{Microbiome} 7 (1): 79.
\url{https://doi.org/10.1186/s40168-019-0694-6}.

\bibitem[\citeproctext]{ref-sanchez_cutaneous_2017}
Sanchez, Eugenia, Molly C Bletz, Laura Duntsch, et al. 2017.
{``Cutaneous Bacterial Communities of a Poisonous Salamander: A
Perspective from Life Stages, Body Parts and Environmental
Conditions.''} \emph{Microbial Ecology} 73 (2): 455--65.
\url{https://doi.org/10.1007/s00248-016-0863-0}.

\bibitem[\citeproctext]{ref-santos_skin_2023}
Santos, Bárbara, Filipa MS Martins, Joana Sabino-Pinto, Fulvio Licata,
and Angelica Crottini. 2023. {``Skin and Gut Microbiomes of Tadpoles
Vary Differently with Host and Water Environment: A Short-Term
Experiment Using {16S} Metabarcoding.''} \emph{Scientific Reports} 13
(1): 16321. \url{https://doi.org/10.1038/s41598-023-43340-2}.

\bibitem[\citeproctext]{ref-scheffers_effects_2012}
Scheffers, Brett R, and Cynthia A Paszkowski. 2012. {``The Effects of
Urbanization on {North} {American} Amphibian Species: {Identifying} New
Directions for Urban Conservation.''} \emph{Urban Ecosystems} 15:
133--47. \url{https://doi.org/10.1007/s11252-011-0199-y}.

\bibitem[\citeproctext]{ref-shaffer_dram_2020}
{Shaffer, Michael, Mikayla A Borton, Bridget B McGivern, et al.} 2020.
{``{DRAM} for Distilling Microbial Metabolism to Automate the Curation
of Microbiome Function.''} \emph{Nucleic Acids Research} 48 (16):
8883--900. \url{https://doi.org/10.1093/nar/gkaa621}.

\bibitem[\citeproctext]{ref-silverman_measuring_2019}
Silverman, Justin D, Rachael J Bloom, Sharon Jiang, Heather K Durand,
Sayan Mukherjee, and Lawrence A David. 2019. {``Measuring and Mitigating
{PCR} Bias in Microbiome Data.''} \emph{BioRxiv}, 604025.
\url{https://doi.org/10.1101/604025}.

\bibitem[\citeproctext]{ref-smith_skin_2018}
Smith, Hannah Keely, Frank Pasmans, Maarten Dhaenens, et al. 2018.
{``Skin Mucosome Activity as an Indicator of {Batrachochytrium}
Salamandrivorans Susceptibility in Salamanders.''} \emph{PLoS One} 13
(7): e0199295. \url{https://doi.org/10.1371/journal.pone.0199295}.

\bibitem[\citeproctext]{ref-soto-pozos_imprints_2025}
Soto-Pozos, Ángel F, Eria A Rebollar, Sean M Rovito, and Gabriela
Parra-Olea. 2025. {``Imprints of {Land} {Use} {History} on the
{Cutaneous} {Microbiota} of {Mexican} {Cloud} {Forest} {Salamanders}.''}
\emph{Microbial Ecology}, ahead of print.
\url{https://doi.org/10.1007/s00248-025-02671-5}.

\bibitem[\citeproctext]{ref-su_neglected_2023}
Su, Rui, Sheng Zhang, Xueying Zhang, Shengnan Wang, and Wenya Zhang.
2023. {``Neglected Skin-Associated Microbial Communities: A Unique
Immune Defense Strategy of {Bufo} Raddei Under Environmental Heavy Metal
Pollution.''} \emph{Environmental Science and Pollution Research} 30
(9): 22330--42. \url{https://doi.org/10.1007/s11356-022-23803-1}.

\bibitem[\citeproctext]{ref-sun_slow_2012}
Sun, Cheng, José R López Arriaza, and Rachel Lockridge Mueller. 2012.
{``Slow {DNA} Loss in the Gigantic Genomes of Salamanders.''}
\emph{Genome Biology and Evolution} 4 (12): 1340--48.
\url{https://doi.org/10.1093/gbe/evs103}.

\bibitem[\citeproctext]{ref-sun_inference-based_2020}
Sun, Shan, Roshonda B Jones, and Anthony A Fodor. 2020.
{``Inference-Based Accuracy of Metagenome Prediction Tools Varies Across
Sample Types and Functional Categories.''} \emph{Microbiome} 8: 1--9.
\url{https://doi.org/10.1186/s40168-020-00815-y}.

\bibitem[\citeproctext]{ref-tattersall_skin_2007}
Tattersall, G. 2007. {``Skin Breathing in Amphibians.''} In
\emph{Endothelial Biomedicine: A Comprehensive Reference}.

\bibitem[\citeproctext]{ref-toledo_cutaneous_1993}
Toledo, RC, and C Jared. 1993. {``Cutaneous Adaptations to Water Balance
in Amphibians.''} \emph{Comparative Biochemistry and Physiology Part A:
Physiology} 105 (4): 593--608.
\url{https://doi.org/10.1016/0300-9629(93)90259-7}.

\bibitem[\citeproctext]{ref-tornabene_energy-related_2023}
Tornabene, Brian J., Kelly L. Smalling, Carrie E. Givens, Emily B. Oja,
and Blake R. Hossack. 2023. {``Energy-Related Wastewater Contamination
Alters Microbial Communities of Sediment, Water, and Amphibian Skin.''}
\emph{Science of The Total Environment} 880 (July): 163160.
\url{https://doi.org/10.1016/j.scitotenv.2023.163160}.

\bibitem[\citeproctext]{ref-touceda-suarez_differences_2025}
Touceda-Suárez, María, Alise J Ponsero, and Albert Barberán. 2025.
{``Differences in the Genomic Potential of Soil Bacterial and Viral
Communities Between Urban Greenspaces and Natural Arid Soils.''}
\emph{Applied and Environmental Microbiology} 91 (8): e02124--24.
\url{https://doi.org/10.1128/aem.02124-24}.

\bibitem[\citeproctext]{ref-uritskiy_metawrapflexible_2018}
Uritskiy, Gherman V, Jocelyne DiRuggiero, and James Taylor. 2018.
{``{MetaWRAP}---a Flexible Pipeline for Genome-Resolved Metagenomic Data
Analysis.''} \emph{Microbiome} 6: 1--13.
\url{https://doi.org/10.1186/s40168-018-0541-1}.

\bibitem[\citeproctext]{ref-velo-anton_salamandra_2015}
Velo-Antón, Guillermo, and David Buckley. 2015. {``Salamandra
Común--{Salamandra} Salamandra ({Linnaeus}, 1758).''} In
\emph{Enciclopedia {Virtual} de Los {Vertebrados} {Españoles}}. Museo
Nacional de Ciencias Naturales, Madrid.

\bibitem[\citeproctext]{ref-wagner_habitat_2020}
Wagner, Norman, Wiebke Harms, Florian Hildebrandt, et al. 2020. {``Do
Habitat Preferences of {European} Fire Salamander ({Salamandra}
Salamandra) Larvae Differ Among Landscapes? {A} Case Study from
{Western} {Germany}.''} \emph{Salamandra} 56 (3).

\bibitem[\citeproctext]{ref-walke_seasonal_2021}
Walke, Jenifer B, Matthew H Becker, Arianna Krinos, et al. 2021.
{``Seasonal Changes and the Unexpected Impact of Environmental
Disturbance on Skin Bacteria of Individual Amphibians in a Natural
Habitat.''} \emph{FEMS Microbiology Ecology} 97 (3): fiaa248.
\url{https://doi.org/10.1093/femsec/fiaa248}.

\bibitem[\citeproctext]{ref-walke_amphibian_2014}
Walke, Jenifer B, Matthew H Becker, Stephen C Loftus, et al. 2014.
{``Amphibian Skin May Select for Rare Environmental Microbes.''}
\emph{The ISME Journal} 8 (11): 2207--17.
\url{https://doi.org/10.1038/ismej.2014.77}.

\bibitem[\citeproctext]{ref-weinstock_genomic_2012}
Weinstock, George M. 2012. {``Genomic Approaches to Studying the Human
Microbiota.''} \emph{Nature} 489 (7415): 250--56.
\url{https://doi.org/10.1038/nature11553}.

\bibitem[\citeproctext]{ref-wheeler_effects_2025}
Wheeler, Sara J, Manning DelCogliano, Brennan Hare, et al. 2025.
{``Effects of Urbanization on the Microbial Fungal Communities from
Water, Sediment, and Amphibian Hosts Across a Rural-to-Urban Waterway in
{Worcester}, {Massachusetts}.''} \emph{bioRxiv}, 2025--04.
\url{https://doi.org/10.1101/2025.04.01.646624}.

\bibitem[\citeproctext]{ref-woodhams_adaptive_2023}
Woodhams, Douglas C., Julia McCartney, Jenifer B. Walke, and Ross
Whetstone. 2023. {``The Adaptive Microbiome Hypothesis and Immune
Interactions in Amphibian Mucus.''} \emph{Developmental \& Comparative
Immunology} 145 (August): 104690.
\url{https://doi.org/10.1016/j.dci.2023.104690}.

\bibitem[\citeproctext]{ref-wu_maxbin_2016}
Wu, Yu-Wei, Blake A Simmons, and Steven W Singer. 2016. {``{MaxBin} 2.0:
An Automated Binning Algorithm to Recover Genomes from Multiple
Metagenomic Datasets.''} \emph{Bioinformatics} 32 (4): 605--7.
\url{https://doi.org/10.1093/bioinformatics/btv638}.

\bibitem[\citeproctext]{ref-yang_global_2022}
Yang, Lu, Shuqing Zhao, and Shuguang Liu. 2022. {``A Global Analysis of
Urbanization Effects on Amphibian Richness: {Patterns} and Drivers.''}
\emph{Global Environmental Change} 73: 102476.
\url{https://doi.org/10.1016/j.gloenvcha.2022.102476}.

\bibitem[\citeproctext]{ref-yang_comparative_2024}
Yang, Peikui, Yuzhong Zheng, Xianghui Zou, Yanjie Sun, and Yaqun Liu.
2024. {``Comparative Transcriptomic Analysis of Gene Expression Profiles
in the Liver and Spleen of {American} Bullfrog ({Lithobates}
Catesbeianus) in Response to {Citrobacter} Freundii Infection.''}
\emph{Journal of the World Aquaculture Society} 55 (1): 353--72.
\url{https://doi.org/10.1111/jwas.12999}.

\bibitem[\citeproctext]{ref-zhang_phylogeny_2008}
Zhang, Peng, Theodore J Papenfuss, Marvalee H Wake, Lianghu Qu, and
David B Wake. 2008. {``Phylogeny and Biogeography of the Family
{Salamandridae} ({Amphibia}: {Caudata}) Inferred from Complete
Mitochondrial Genomes.''} \emph{Molecular Phylogenetics and Evolution}
49 (2): 586--97. \url{https://doi.org/10.1016/j.ympev.2008.08.020}.

\bibitem[\citeproctext]{ref-zheng_urbanization_2025}
Zheng, Bangxiao, Nan Hui, Ari Jumpponen, et al. 2025. {``Urbanization
Leads to Asynchronous Homogenization of Soil Microbial Communities
Across Biomes.''} \emph{Environmental Science and Ecotechnology} 25:
100547. \url{https://doi.org/10.1016/j.ese.2025.100547}.

\bibitem[\citeproctext]{ref-zhou_urbanization_2023}
Zhou, Jin, Ziyan Liao, Zhidong Liu, Xuecheng Guo, Wenyan Zhang, and
Youhua Chen. 2023. {``Urbanization Increases Stochasticity and Reduces
the Ecological Stability of Microbial Communities in Amphibian Hosts.''}
\emph{Frontiers in Microbiology} 13: 1108662.
\url{https://doi.org/10.3389/fmicb.2022.1108662}.

\bibitem[\citeproctext]{ref-zhu_skin_2024}
Zhu, De-quan, Wen-jing Dong, Xin-zhou Long, et al. 2024. {``Skin Ulcers
and Microbiota in {Rana} Dybowskii: Uncovering the Role of the Gut-Skin
Axis in Amphibian Health.''} \emph{Aquaculture} 585: 740724.
\url{https://doi.org/10.1016/j.aquaculture.2024.740724}.

\bibitem[\citeproctext]{ref-ziesemer_intrinsic_2015}
{Ziesemer, Kirsten A, Allison E Mann, Krithivasan Sankaranarayanan, et
al.} 2015. {``Intrinsic Challenges in Ancient Microbiome Reconstruction
Using {16S} {rRNA} Gene Amplification.''} \emph{Scientific Reports} 5
(1): 16498. \url{https://doi.org/10.1038/srep16498}.

\end{CSLReferences}

\end{document}
